\documentclass[11pt]{article}

\usepackage[margin=1in]{geometry}
\usepackage[utf8]{inputenc}
\usepackage[T1]{fontenc}
\usepackage{amsmath,amssymb}
\usepackage{booktabs}
\usepackage{array}
\usepackage{graphicx}
\usepackage{xcolor}
\usepackage{enumitem}
\usepackage{caption}
\usepackage{longtable}
\usepackage{float}
\usepackage{algorithm}
\usepackage{algpseudocode}
\usepackage[hidelinks,colorlinks=false]{hyperref}
\usepackage[noabbrev,capitalise]{cleveref}

\setlength{\parindent}{0pt}
\setlength{\parskip}{6pt}
\emergencystretch=3em

\definecolor{cIndigo}{HTML}{3A4CC4}
\definecolor{cAmber}{HTML}{96590A}
\definecolor{cInk2}{HTML}{525A6C}
\definecolor{cPanel}{HTML}{F6F7F9}
\definecolor{cRule}{HTML}{BFC6D3}

\newcommand{\dgraph}{d_{\mathrm{graph}}}
\newcommand{\zv}{\mathbf{z}}

\title{\vspace{-2em}\textbf{Latent Graph Shaping}\\[0.3em]
\large A reward-free, calibrated distance from LeWM's latent,\\
deployed as a policy-invariant potential for offline goal-conditioned RL}
\author{Experiment proposal}
\date{28 August 2026}

\begin{document}
\maketitle
\vspace{-1.5em}

\section{What this proposes, in one paragraph}

Offline RL's distinctive capability over behaviour cloning is \emph{stitching}: recombining
sub-trajectories $A\!\to\!B$ and $B\!\to\!C$ into $A\!\to\!B\!\to\!C$. Standard algorithms
obtain this only indirectly, through many steps of TD bootstrapping. We propose to compute it
directly and in one pass: encode dataset states with a frozen pretrained LeWM encoder, build a
graph whose edges are (i) real recorded transitions and (ii) cross-trajectory
\emph{identifications} between states the latent says are the same, and run Dijkstra. The
resulting distance $\dgraph(s,g)$ is then supplied to an off-the-shelf offline algorithm as a
potential-based shaping term, which \emph{cannot change the optimal policy} (Ng, Harada \&
Russell 1999) --- so a bad graph costs speed, never correctness. Nothing is added to the RL
inner loop. No reward labels are used anywhere in the construction.

The identification threshold is not hand-picked: it comes from the same $\chi^2$
one-step-displacement null already used elsewhere in this project, with the correlation
$\hat\rho$ estimated from consecutive dataset frames.

\textbf{Status.} The gate experiment (B0, \cref{sec:b0-results}) is done and passed decisively
on Two-Room with the real pretrained checkpoint. B1, B2, B3, and B4 are also done
(\cref{sec:b1-results,sec:b2-results,sec:b3-results,sec:b4-results}), including a live-rollout
confirmation that the B3/B4 value-ranking win converts into real task-success
(\cref{sec:actor-rollout}) and a B4 extension confirming that B1's required edge-count cap is a
net improvement, not a cost, at the RL-training level (\cref{sec:b4ext-results}). A live-rollout
confirmation of that extension (\cref{sec:b4ext-actor-rollout}) qualifies the proxy-metric result:
$k$-capping over baseline holds up, but $k$-capping beating the \emph{full} graph specifically was
a proxy-metric artifact. \Cref{sec:pusht} then transfers the whole approach to a second
environment, Push-T: the gate check (\cref{sec:b5-results}) is done, and passes via the kill
criterion's named fallback rather than the primary calibrated path; the RL-training pipeline and
its live-rollout confirmation are also done, across several rounds of diagnosis and fixes
summarised in that section.

\section{Method}
\label{sec:method}

\subsection{Construction}

\begin{algorithm}[H]
\caption{Latent graph construction}
\label{alg:construction}
\begin{algorithmic}[1]
\Require dataset trajectories with frames $\{o_i\}$; frozen pretrained encoder $\mathrm{enc}(\cdot)$; calibration quantile $q$
\Ensure a weighted graph $G$ over encoded landmarks, queryable as $\dgraph(s,g)$
\State $\zv_i \gets \mathrm{enc}(o_i)$ for every frame \Comment{$d=192$; one forward pass, cached once; subsample to landmarks for large datasets}
\State $\hat\rho \gets \dfrac{1}{d}\,\mathbb{E}_{\mathcal D}\bigl[\zv_t^\top \zv_{t+1}\bigr]$ \Comment{one-step displacement correlation, estimated from the data}
\State $\varepsilon^2 \gets 2(1-\hat\rho)\,F^{-1}_{\chi^2_d}(q)$ \Comment{small-$q$ (left) tail of the one-step-displacement null}
\For{each recorded trajectory}
    \For{each consecutive pair of frames $(i, i{+}1)$}
        \State add directed \emph{transition} edge $(i,i{+}1)$, weight $1$ \Comment{exact and free --- it is the data}
    \EndFor
\EndFor
\For{each cross-trajectory pair $(i,j)$}
    \If{$\|\zv_i - \zv_j\|^2 < \varepsilon^2$}
        \State add undirected \emph{identification} edge $(i,j)$, weight $1$ \Comment{weight $0$ breaks the method --- \cref{sec:bugs}}
    \EndIf
\EndFor
\Return $\dgraph(s,g) \gets$ Dijkstra shortest path from $s$ to $g$ on $G$ \Comment{milliseconds per query}
\end{algorithmic}
\end{algorithm}

An identification edge admits pairs strictly closer in latent space than a typical single step
($\varepsilon^2 = 114$ vs.\ mean one-step $\|\Delta \zv\|^2 = 158$ on Two-Room) --- weighting them
at $1$ step is a conservative upper bound on their true cost, and requires no further tuning.

\subsection{How it is consumed}

\begin{algorithm}[H]
\caption{Consuming $\dgraph$ in offline RL --- two modes tested, in trial order}
\label{alg:consumption}
\begin{algorithmic}[1]
\Require $\dgraph(s,g)$ from \cref{alg:construction}; goal-conditioned value/critic $V(s,g)$, $Q(s,a,g)$; discount $\gamma$
\Procedure{Mode3-PotentialShaping}{$s,s',g$} \Comment{tried first --- recommended \emph{a priori}}
    \State $\Phi(s) \gets -\dgraph(s,g)$
    \Return $r(s,s') + \gamma\Phi(s') - \Phi(s)$ \Comment{policy-invariant for any $\Phi$ (Ng, Harada \& Russell 1999); B3: loses on 4/5 tiers}
\EndProcedure
\Procedure{Mode2-AuxRegression}{$s,g$} \Comment{tried second}
    \State $\hat v \gets -\bigl(1-\gamma^{\dgraph(s,g)}\bigr)/(1-\gamma)$ \Comment{pseudo-value under pure geometric discounting}
    \Return $\lambda_{\text{aux}}\,\bigl(V(s,g) - \hat v\bigr)^2$ \Comment{added to $V$'s ordinary loss; TD objective itself untouched; B3: wins at every tier tested}
\EndProcedure
\end{algorithmic}
\end{algorithm}

Mode 3 was the point of the design going in: potential-based shaping's policy-invariance guarantee
means a bad graph can only cost speed, never correctness, since nothing in the RL objective
searches against $\Phi$ the way an adversarial optimisation would. That safety property held. What
did not hold was the assumption that safety plus a very accurate $\Phi$ would translate into a
speed-up in practice --- mode 2, tried second, is the one B3 finds actually works.

\begin{center}
\small
\begin{longtable}{@{}p{0.5cm}p{3.5cm}p{5.1cm}p{4.6cm}@{}}
\caption{Experiment ladder.} \\
\toprule
& \textbf{Question} & \textbf{Method} & \textbf{Kill criterion} \\
\midrule
\endfirsthead
\multicolumn{4}{l}{\small\itshape (Experiment ladder, continued)} \\
\toprule
& \textbf{Question} & \textbf{Method} & \textbf{Kill criterion} \\
\midrule
\endhead
\bottomrule
\multicolumn{4}{r}{\small\itshape (continued on next page)} \\
\endfoot
\bottomrule
\endlastfoot
B0 & Does the graph metric beat raw latent distance? & Spearman vs.\ oracle geodesic, Two-Room.
& \textcolor{cAmber}{\textbf{Done --- passed}} (0.95 vs.\ 0.44). \\
\addlinespace
B1 & Is the graph well-behaved, and does it survive scale? & Landmark-count scaling
(50/300/1000 eps); successor-consistency filter vs.\ false-edge rate; connectivity/fragmentation
vs.\ $q$; degree distribution. & \textcolor{cAmber}{\textbf{Done --- mixed.}} Quality survives
scale; edge count does not (\cref{sec:b1-results}). \\
\addlinespace
B2 & Is this about \emph{LeWM}, or about any latent? & Swap the encoder, same pipeline, same
$q$: random-init ViT, raw pixels, PCA-192 on the pooled pixels (done); DINO-WM, PLDM, LeJEPA
(blocked, dead Drive link, no alternative source). & \textcolor{cAmber}{\textbf{Partially done ---
did not fire.}} Training matters; whether it's \emph{SIGReg} specifically remains open
(\cref{sec:b2-results}). \\
\addlinespace
B3 & \textbf{Does it actually help an RL agent?} & GCIQL / GCIVL $\pm$ graph-potential shaping.
Metrics: steps-to-threshold (the speed claim) and final success rate. Across data-quality tiers
(expert / mixed / subsampled). & \textcolor{cAmber}{\textbf{Scoped pilot done --- reward shaping
mostly did not fire, but a fix does.}} See \cref{sec:b3-results,sec:actor-rollout}. \\
\addlinespace
B4 & Is the potential doing the work, or just any potential? & Ablate the distance source fed
to mode 2 (auxiliary regression, the mechanism B3 found actually works): none (baseline) /
raw Euclidean / graph geodesic / true-geodesic oracle. All six tiers, 3 seeds. &
\textcolor{cAmber}{\textbf{Done --- kill criterion did not fire.}} Graph beats Euclidean at
every tier, margin growing with data (\cref{sec:b4-results}). Extension: $k$-capped graphs
(k4/k8) beat even the uncapped graph at every tier on the Spearman proxy
(\cref{sec:b4ext-results}); on real rollout success they instead match the full graph and still
clearly beat baseline (\cref{sec:b4ext-actor-rollout}). \\
\addlinespace
Push-T & Does the whole approach transfer to a second environment? & Repeat B0's gate check and
B3's data-tier sweep on Push-T, plus a live-rollout confirmation; $q$ held fixed at the
Two-Room value initially. & \textcolor{cAmber}{\textbf{Done --- transfers, with a per-environment
calibration fix and several diagnosed fixes along the way.}} See \cref{sec:pusht} for the full
account (gate check, RL-training, live-rollout confirmation, and a stitching-specific showcase). \\
\end{longtable}
\end{center}

\section{Experiments}

B0, B1, B2, and B4 are complete; B3 is complete in its scoped form (GCIQL only, a value-ranking
proxy metric, plus a later live-rollout confirmation at the real task-success level). The table
above summarises all five, all on Two-Room; each subsection below walks through one experiment's
methodology, results, and a brief takeaway, in that order. \Cref{sec:pusht} then asks the same
questions of a second environment, Push-T, as its own section.

\subsection{B0: does the graph metric beat raw latent distance?}
\label{sec:b0-results}

\textbf{Methodology.}

\begin{algorithm}[H]
\caption{B0 methodology}
\begin{algorithmic}[1]
\State \textbf{Checkpoint:} real pretrained \texttt{quentinll/lewm-tworooms} \Comment{verified not collapsed: off-diag.\ cosine $\approx 0.03$ vs.\ $\approx 0.9999$ for the local epoch-1 checkpoint}
\State \textbf{Data:} 300 episodes $\to$ 27{,}675 encoded landmarks, $\hat\rho = 0.584$
\State \textbf{Graph:} build via \cref{alg:construction}, $q = 10^{-3}$
\State \textbf{Ground truth:} grid-Dijkstra oracle over recorded $(x,y)$ positions \Comment{wall geometry read empirically off trajectories, not env source defaults}
\State \textbf{Eval set:} 3{,}999 held-out query pairs
\State \textbf{Metric:} Spearman($\dgraph$, true distance) vs.\ Spearman($\|\zv_i-\zv_j\|$, true distance)
\State \textbf{Robustness check:} repeat for $q \in \{0.5,\,0.05,\,0.01,\,10^{-3},\,10^{-6},\,10^{-9}\}$, edge weight held fixed
\end{algorithmic}
\end{algorithm}

\textbf{Results.}

\begin{table}[h]
\centering
\caption{B0 result. Spearman correlation with true wall-respecting distance, at the principled
setting (identification weight $=1$, $q=10^{-3}$). All $p \approx 0$.}
\begin{tabular}{lcc}
\toprule
& raw latent Euclidean & graph geodesic \\
\midrule
overall & 0.44 & \textbf{0.95} \\
same-room & 0.33 & \textbf{0.93} \\
cross-room (through the door) & 0.66 & \textbf{0.94} \\
\bottomrule
\end{tabular}
\end{table}

\begin{table}[h]
\centering
\caption{Sensitivity to the calibration quantile $q$. ``False edge'' = identification edge whose
true geodesic separation exceeds two environment steps.}
\begin{tabular}{lccccc}
\toprule
$q$ & $\varepsilon^2$ & avg.\ degree & overall & cross-room & false edges \\
\midrule
$0.5$      & 159.2 & 275 & 0.839 & 0.709 & 6.8\% \\
$0.05$     & 134.0 & 223 & 0.900 & 0.829 & 4.6\% \\
$0.01$     & 124.3 & 205 & 0.934 & 0.905 & 4.5\% \\
$10^{-3}$  & 114.1 & 186 & 0.946 & 0.930 & 3.4\% \\
$10^{-6}$  &  93.9 & 151 & 0.966 & 0.981 & 2.5\% \\
$10^{-9}$  &  80.5 & 128 & 0.972 & 0.986 & 2.2\% \\
\bottomrule
\end{tabular}
\end{table}

Performance is \emph{not} flat in $q$ --- it improves monotonically as $q$ tightens, tracking the
false-edge rate. Every value across five decades ($q \le 10^{-2}$) still gives overall
$\ge 0.93$.

\label{sec:bugs}
Two implementation bugs were caught during this process by measuring things that could have been
assumed, and both would have silently produced a wrong negative result: (i) weighting
identification edges at $0$ (``same state, free hop'') makes the graph \emph{worse} than raw
Euclidean distance (cross-room $0.21$ vs.\ $0.66$), because Dijkstra chains many free
$\sim$4\,px hops into large illegitimate shortcuts --- charging one step per hop fixes it; (ii)
the $\chi^2$ tail direction was initially inverted ($F^{-1}(1-\alpha)$ instead of $F^{-1}(q)$),
which makes the threshold \emph{looser} as the nominal level tightens and inflated average degree
from 32 to 27{,}580 as $\alpha$ shrank --- caught by plotting edge count against the knob and
noticing the trend ran backwards.

\textbf{Takeaways.}
\begin{itemize}[leftmargin=1.4em,itemsep=2pt]
\item Graph geodesic decisively beats raw latent distance (0.95 vs.\ 0.44 overall Spearman),
including on same-room pairs where there is no topological reason for it to help --- suggesting
it also denoises the latent, not merely routes around the wall.
\item The result is robust, not flat: quality tracks $q$ monotonically, but the defensible claim
is ``robust across five decades'' rather than ``a single untuned $\alpha$ works.''
\item Two edge-construction bugs (zero-weight identifications; inverted $\chi^2$ tail) were
caught by testing before they could produce a silent wrong negative --- both fixed.
\end{itemize}

\subsection{B1: is the graph well-behaved, and does it survive scale?}
\label{sec:b1-results}

\textbf{Methodology.}

\begin{algorithm}[H]
\caption{B1 methodology}
\begin{algorithmic}[1]
\State \textbf{Base setup:} same checkpoint and pipeline as B0, $q=10^{-3}$ fixed, identification weight fixed at $1$
\State \textbf{Landmark scaling:} repeat \cref{alg:construction} at $N \in \{50, 300, 1000\}$ episodes; record avg.\ degree, edge count, overall/cross-room Spearman
\For{each identification edge $(i,j)$ with $\cos(a_i,a_j) > 0.7$}
    \State \textbf{successor-consistency test:} also require successors $(i{+}1,\,j{+}1)$ pass the identification test; drop $(i,j)$ if they don't
\EndFor
\State \textbf{Connectivity:} sweep $q$ from $0.5$ down to $10^{-9}$; record fragmentation and cross-room reachability
\State \textbf{Degree distribution:} at $q=10^{-3}$, report mean / p99 / max degree
\end{algorithmic}
\end{algorithm}

\textbf{Results.}

\begin{table}[h]
\centering
\caption{Scaling with landmark budget. Quality degrades gently, edge count does not.}
\begin{tabular}{lccccc}
\toprule
episodes & landmarks & avg.\ degree & overall & cross-room & id.\ edges \\
\midrule
50   &  4{,}595 &  32 & 0.966 & 0.984 & 74{,}326 \\
300  & 27{,}675 & 186 & 0.946 & 0.930 & 2{,}578{,}506 \\
1000 & 92{,}552 & 578 & 0.930 & 0.887 & 26{,}744{,}539 \\
\bottomrule
\end{tabular}
\end{table}

Overall Spearman only drifts 0.966 $\to$ 0.930 across a 20$\times$ landmark range, staying far
above raw Euclidean (0.44--0.49) throughout. Average degree, however, scales roughly linearly
with landmark count, so total identification edges scale as
$n_{\mathrm{id}} \sim n_{\mathrm{landmarks}}^{1.96}$ --- essentially quadratic, extrapolating to
$\sim 1.5 \times 10^{10}$ edges at Push-T's 2.3M frames.

Successor-consistency filtering: only \textbf{25.6\%} of identification edges were even testable
under the action-similarity condition, because most identified pairs have near-uncorrelated
actions (mean cosine $\approx 0.02$) --- Two-Room's target is never rendered into the pixel
observation, so two episodes with the agent at the same spot but different (invisible) goals head
in unrelated directions. Of the testable 25.6\%, 33\% failed and were dropped (8.4\% of all
identification edges). Net effect: false-edge rate 3.7\% $\to$ 3.8\% (unchanged), downstream
Spearman 0.9455 $\to$ 0.9456 (unchanged).

Connectivity is full at every $q$ tested, $0.5$ down to $10^{-9}$ --- cross-room reachability
stays at 100\% throughout, though this is not really identification edges' doing: 446/500 sampled
episodes already cross the door within their own single recorded trajectory. Degree distribution
at $q=10^{-3}$ is clean: mean 186, p99 423, max 657 --- a smooth tail, zero pathological hubs.

\textbf{Takeaways.}
\begin{itemize}[leftmargin=1.4em,itemsep=2pt]
\item Quality does not collapse with scale --- the originally-written B1 kill criterion did not
fire.
\item But edge count explodes ($\sim n^{1.96}$), extrapolating to infeasibility at Push-T scale.
\textbf{This is the real kill criterion, and it already fired}: a top-$k$-per-node cap is now a
prerequisite for the Push-T transfer work (\cref{sec:pusht}), not an optional efficiency tweak.
\item Successor-consistency filtering does not help on Two-Room: only 25.6\% of edges are even
testable (the invisible goal starves action-similarity of signal), and the net effect on
false-edge rate and Spearman is negligible. May work better where the observation determines the
policy more fully --- should not be assumed for Push-T without checking.
\item Identification edges are doing metric work (B0's Spearman gain) on Two-Room, not
connectivity work --- transition edges alone already carry most of the room-crossing
connectivity here.
\end{itemize}

\subsection{B2: is this about \emph{LeWM}, or about any latent?}
\label{sec:b2-results}

\textbf{Methodology.}

\begin{algorithm}[H]
\caption{B2 methodology}
\begin{algorithmic}[1]
\State \textbf{Fixed:} same 300 episodes, same frames, same pipeline ($q=10^{-3}$, identification weight $1$)
\For{$\mathrm{enc} \in \{$LeWM (pretrained), random-init ViT-tiny, raw pixels ($16{\times}16$ pooled), PCA-192 on the pooled pixels$\}$}
    \State re-run \cref{alg:construction} and B0's evaluation with this encoder in place of LeWM
\EndFor
\State \textbf{PCA basis:} 192 components, fit via IncrementalPCA over the \emph{entire} dataset (920{,}809 frames, all 10{,}000 episodes) --- not just the 300-episode landmark subset it is then evaluated on
\If{identification-edge count $> 10^7$ at 300 episodes}
    \State fall back to the same four encoders on a 10-episode subset (same episodes, same reused encodings, no re-encoding) \Comment{avg.\ degree scales $\sim$linearly with landmark count here, unlike LeWM (B1), so a smaller subset is tractable}
\EndIf
\State \textbf{Scope note:} random-init, raw pixels, and PCA test the \emph{weak} form (does training/dimensionality reduction alone matter), not the \emph{sharp} form (is it SIGReg specifically) \Comment{DINO-WM/PLDM/LeJEPA baselines, which would test the sharp form, blocked by a dead Drive link, confirmed dead with no alternative source}
\end{algorithmic}
\end{algorithm}

\textbf{Results.}

\begin{table}[h]
\centering
\caption{B2 encoder swap. Squared-distance scale and graph Spearman vary enormously by encoder.}
\begin{tabular}{lccccc}
\toprule
encoder & $d$ & $\hat\rho$ (adjacent) & overall (graph) & same-room & cross-room \\
\midrule
LeWM (pretrained)        & 192 & 0.584 & \textbf{0.946} & \textbf{0.932} & \textbf{0.930} \\
random-init ViT-tiny     & 192 & 1.000 & $-0.015$        & 0.087          & $-0.192$ \\
raw pixels (16$\times$16 pooled) & 768 & 0.810 & \multicolumn{3}{c}{intractable (see below)} \\
PCA-192 (on the pooled pixels) & 192 & 0.005 & \multicolumn{3}{c}{intractable (see below)} \\
\bottomrule
\end{tabular}
\end{table}

Random-init ViT: $\hat\rho \approx 1.000$ means adjacent frames map to nearly identical vectors
regardless of content --- a known failure mode of deep random (untrained) networks. Both raw
Euclidean and graph-geodesic distance score at or below zero correlation with true distance.

Raw pixels: 381{,}644{,}170 identification edges at $q=10^{-3}$ --- more than LeWM (2.6M) and
random-ViT (101M) combined --- would need $\sim$6GB just to materialize the edge-index arrays,
which crashes (a hard native segfault, not a catchable Python exception) on this machine. Sweeping
$q$ down to $10^{-80}$ changed \emph{nothing}: identical edge count every time. Random
cross-episode pairs in 16$\times$16-pooled pixel space have a median squared distance of $2.24$
--- the fixed wall/door background dominates the 768-dimensional vector so completely that
``adjacent'' and ``unrelated'' pairs are barely distinguishable in Euclidean terms at any
threshold.

PCA-192: motivated by exactly that diagnosis --- PCA's entire job is finding the directions that
actually vary and discarding near-constant ones (the fixed background), so it looked like a
near-free candidate fix. It is not: \textbf{381{,}644{,}170} identification edges ---
indistinguishable from raw pixels' count to six significant figures --- and a median cross-episode
squared distance of $2.2376$, essentially identical to raw pixels' $2.2380$. $\hat\rho$ drops to
$0.005$ (versus raw pixels' $0.810$), meaning adjacent frames are no longer even weakly
correlated once projected. In hindsight this is the expected result, not a surprising one: PCA is
a \emph{linear}, near-lossless map (192 of 768 components already capture $99.97\%$ of variance)
--- it re-expresses the same raw-pixel geometry in a rotated, whitened basis, but introduces no
new invariance, so a separation problem that is already present in that geometry survives the
projection essentially unchanged.

Edge count for raw pixels/PCA scales very differently from LeWM's with landmark count: B1 found
LeWM's average degree stays roughly \emph{constant} as data grows (32 at 50 episodes, 186 at 300,
578 at 1000), so total edges scale sub-quadratically; raw pixels and PCA instead have avg.\ degree
scaling roughly \emph{linearly} with landmark count (checked directly: 893 at 10 episodes, 9{,}289
at 100, 18{,}544 at 200), so total edges scale roughly quadratically and 10 episodes is already
the largest tractable subset on this machine. Re-running all four encoders on that same 10-episode
subset (same episodes throughout, reusing the already-computed encodings) gives a real,
apples-to-apples number instead of just ``intractable'':

\begin{table}[h]
\centering
\caption{Small-scale (10 episodes, 992 frames) comparison, same episode subset for all four
encoders.}
\begin{tabular}{lccccc}
\toprule
encoder & $d$ & overall (eucl.) & overall (graph) & same-room (graph) & cross-room (graph) \\
\midrule
LeWM (pretrained)     & 192 & 0.438 & \textbf{0.938} & \textbf{0.944} & \textbf{0.932} \\
random-init ViT-tiny  & 192 & $-0.108$ & $-0.108$ & 0.064 & $-0.424$ \\
raw pixels            & 768 & 0.619 & $-0.118$ & $-0.173$ & $-0.001$ \\
PCA-192               & 192 & 0.618 & $-0.118$ & $-0.173$ & $-0.001$ \\
\bottomrule
\end{tabular}
\end{table}

This is sharper than ``intractable,'' not just a fallback: at a scale where the graph geodesic
\emph{can} be computed, raw pixels and PCA's graph distance is not merely uncorrelated with truth,
it is \emph{worse than doing nothing} --- raw Euclidean distance in the same space scores $0.619$
/ $0.618$, but routing through the graph drags it down to $-0.118$ for both, again matching to
three decimal places. With avg.\ degree $893$ at only $992$ nodes, the identification edges are
dense enough that Dijkstra chains illegitimate shortcuts almost everywhere --- the same mechanism
already diagnosed in B0's zero-weight-edge bug (\cref{sec:bugs}), except here it is not a
weighting mistake, just too many genuinely-under-threshold-but-not-actually-nearby pairs for the
threshold to mean anything.

\textbf{Takeaways.}
\begin{itemize}[leftmargin=1.4em,itemsep=2pt]
\item Random-init ViT is collapsed, not merely worse ($\hat\rho\approx1$, Spearman $\le 0$) ---
evidence \emph{for} the trained representation doing real work, not just architecture. The sharp
kill criterion did not fire.
\item Raw pixels and PCA are intractable at 300 episodes but, checked at the largest tractable
subset (10 episodes), the graph is not just unhelpful --- it is actively \emph{worse than raw
Euclidean distance} in the same space, for both, to three decimal places.
\item PCA does not fix raw pixels' failure: it reproduces its edge count, squared-distance
distribution, and small-scale graph Spearman almost exactly, because a linear, variance-preserving
projection cannot introduce the invariances (to the fixed background, in particular) that the
problem actually needs --- evidence that the failure is about representation \emph{content}, not
merely dimensionality.
\item Scope caveat: none of random-init, raw pixels, or PCA test whether it's \emph{SIGReg}
specifically versus another trained objective --- that comparison remains blocked (dead Drive
link, confirmed with no alternative source) and open.
\end{itemize}

\subsection{B3: does it actually help an RL agent?}
\label{sec:b3-results}

\textbf{Methodology.}

\begin{algorithm}[H]
\caption{B3 methodology: setup and the initial shaping test}
\begin{algorithmic}[1]
\State \textbf{Algorithm:} GCIQL only (not GCIVL); expectile loss reused verbatim from \texttt{stable\_worldmodel.wm.gcrl.ExpectileLoss}
\State \textbf{Split:} 80/20 train/test episodes per tier \Comment{graph built over \emph{all} episodes, since $\Phi$ is a fixed deterministic function of the frozen encoder, not fit to minimise eval error}
\State \textbf{HER:} tuples drawn only from train episodes
\State \textbf{Eval sets:} 1{,}000 fixed train pairs, 1{,}000 fixed held-out test pairs; true distances precomputed once
\State \textbf{Tiers:} expert\_$\{10,25,50,100\}$ (real episodes only); mixed $=$ 50 real $+$ 50 (later 220) live noisy rollouts, \texttt{ExpertPolicy(action\_noise{=}0.4, action\_repeat\_prob{=}0.15)}
\State \textbf{Conditions:} baseline ($-1/0$ TD, untouched) vs.\ baseline $+$ \Call{Mode3-PotentialShaping}{} (\cref{alg:consumption})
\State \textbf{Training:} 50{,}000 steps, evaluate every 50 steps; checkpoint every eval step, resume on interrupt
\State \textbf{Metric:} peak held-out-test Spearman, not final-step value \Comment{every condition on every tier peaks early then degrades from a shared TD/overfitting instability}
\end{algorithmic}
\end{algorithm}

\begin{algorithm}[H]
\caption{B3 methodology: mechanism diagnosis and fix attempts, once shaping was found to fail}
\begin{algorithmic}[1]
\State \textbf{Graph-pollution check:} compare Spearman($\Phi$, true) on expert-only pairs, with vs.\ without noisy-tier nodes in the graph
\State \textbf{Disconnection check:} connectivity, reachability, and identification-edge degree, split by node origin (expert vs.\ noisy)
\State \textbf{Reward-noise-mismatch check:} mean/std of $\gamma\Phi(s')-\Phi(s)$ on HER tuples, split by origin episode
\State \textbf{Peak-snapshot inspection:} re-train keeping peak-step (not final-step) weights; inspect $V(s,g)$ shape against true-distance bin
\State \textbf{Fix attempt 1:} $\lambda\cdot(\gamma\Phi(s')-\Phi(s))$ for $\lambda \in \{0.1, 0.15, 0.3, 0.5\}$ \Comment{preserves the exact potential-based structure}
\State \textbf{Fix attempt 2:} replace shaping with \Call{Mode2-AuxRegression}{} (\cref{alg:consumption}), $\lambda_{\text{aux}}=0.3$, all six tiers, 3 seeds
\end{algorithmic}
\end{algorithm}

\textbf{Results.}

\begin{table}[h]
\centering
\caption{B3 across data tiers: peak held-out-test Spearman, mean $\pm$ std over 3 seeds.}
\begin{tabular}{lccc}
\toprule
tier & baseline & shaped & shaped $-$ baseline \\
\midrule
expert (10 eps)  & $0.303 \pm 0.013$ & $0.221 \pm 0.018$ & $-0.082$ \\
expert (25 eps)  & $0.197 \pm 0.008$ & $0.145 \pm 0.029$ & $-0.052$ \\
expert (50 eps)  & $0.272 \pm 0.030$ & $0.212 \pm 0.049$ & $-0.061$ \\
expert (100 eps) & $0.306 \pm 0.024$ & $\mathbf{0.328 \pm 0.028}$ & $\mathbf{+0.022}$ \\
mixed (50+50 eps) & $\mathbf{0.340 \pm 0.013}$ & $0.205 \pm 0.012$ & $-0.135$ \\
\bottomrule
\end{tabular}
\end{table}

Shaping (mode 3) wins at exactly one of five tiers: the largest, cleanest one (100 clean expert
episodes, shaped peak higher in under half the steps). Everywhere else it loses, worst on the
mixed-quality tier. Scaling the mixed tier up further (\texttt{mixed\_large}: same 50 real
episodes, 220 noisy episodes instead of 50, $\sim$10{,}400 landmarks) makes this worse, not
better: baseline nearly doubles (0.340 $\to$ 0.575), shaping barely moves (0.205 $\to$ 0.266) ---
the gap widens from $-0.135$ to $-0.309$, the worst of any tier tested.

Mechanism checks: Spearman($\Phi$, true distance) on identical expert-only pairs is the same
(0.969 vs.\ 0.967) whether or not the graph also contains noisy-tier nodes --- not graph
pollution. One connected component, 100\% of noisy-origin nodes reachable, comparable
identification-edge degree by origin (29.5 vs.\ 33.4) --- not disconnection. The shaping
increment $\gamma\Phi(s')-\Phi(s)$ has mean $0.46$, std $0.68$ on real-expert HER tuples vs.\
mean $0.75$, std $0.47$ on noisy-collector tuples: real trajectories track monotonic
graph-distance progress toward HER-sampled goals \emph{less} consistently (91 avg.\ steps/episode
vs.\ $\sim$25 for the direct waypoint-following collector) even though $\Phi$ itself stays
uniformly accurate --- a population-conditional reward-noise mismatch that mode 3 folds into the
reward, where TD bootstraps and compounds it; plain TD's $-1/0$ reward has no such channel.
Peak-snapshot inspection: baseline's peak $V(s,g)$ is nearly flat (range $0.08$) but ranks
decently ($\rho=0.21$); shaped's peak $V$ has $\sim\!10\times$ more range ($0.81$) but is
non-monotonic --- correct near and far, reversed at mid-range distances --- and that reversal is
exactly where it loses ($\rho=0.17$).

Fix attempt 1 ($\lambda$-scaled mode 3): on the mixed tier, $\lambda=0.1\to0.328$,
$0.15\to0.323$, $0.3\to0.311$, $0.5\to0.293$ --- recovers $\sim$90\% of the $\lambda=1$ deficit
($0.205$) but never crosses baseline's $0.340$.

\begin{table}[h]
\centering
\small
\caption{B3 final: peak held-out-test Spearman, mean $\pm$ std over 3 seeds, all six tiers.}
\begin{tabular}{lcccc}
\toprule
tier & baseline & shaped (mode 3) & \textbf{aux.\ reg.\ (mode 2)} & $\Delta$ (mode 2) \\
\midrule
expert\_10   & $0.303 \pm 0.013$ & $0.221 \pm 0.018$ & $0.316 \pm 0.039$ & $+0.013$ \\
expert\_25   & $0.197 \pm 0.008$ & $0.145 \pm 0.029$ & $\mathbf{0.351 \pm 0.028}$ & $\mathbf{+0.154}$ \\
expert\_50   & $0.272 \pm 0.030$ & $0.212 \pm 0.049$ & $\mathbf{0.383 \pm 0.012}$ & $\mathbf{+0.111}$ \\
expert\_100  & $0.306 \pm 0.024$ & $0.328 \pm 0.028$ & $\mathbf{0.646 \pm 0.010}$ & $\mathbf{+0.339}$ \\
mixed        & $0.340 \pm 0.013$ & $0.205 \pm 0.012$ & $\mathbf{0.490 \pm 0.016}$ & $\mathbf{+0.150}$ \\
mixed\_large & $0.575 \pm 0.007$ & $0.266 \pm 0.026$ & $\mathbf{0.908 \pm 0.002}$ & $\mathbf{+0.334}$ \\
\bottomrule
\end{tabular}
\end{table}

Fix attempt 2 (mode 2, auxiliary regression) beats baseline at every tier tested, by a margin
that grows with data --- the precise inverse of mode 3's pattern. Only expert\_10 is a wash (seed
ranges overlap); from expert\_25 up there is no overlap in either direction. On
\texttt{mixed\_large}, mode 2 reaches $0.908$, beating baseline's already-strong $0.575$ by
$+0.334$.

\textbf{Takeaways.}
\begin{itemize}[leftmargin=1.4em,itemsep=2pt]
\item Potential-based shaping (mode 3) wins only at the largest/cleanest tier, loses everywhere
else, and gets \emph{worse} rather than better as more heterogeneous data is added --- inverting
the proposal's own hypothesis that shaping would help most when data is scarce.
\item Root cause: not graph noise ($\Phi$ stays $\ge0.96$ accurate throughout) but a
population-conditional reward-noise mismatch that TD bootstrapping compounds across a trajectory;
plain $-1/0$ TD has no such channel, which is why it alone benefits from added heterogeneous data.
\item Scaling mode 3's magnitude down recovers most, not all, of the deficit --- confirms a
magnitude/variance issue, not a fix worth shipping.
\item The fix that works: stop routing $\Phi$ through the reward at all. Mode 2 (direct,
non-bootstrapped auxiliary regression toward a $\Phi$-derived pseudo-value) beats baseline at
every tier, margin growing with data ($+0.34$ at expert\_100 and mixed\_large).
\item Label efficiency (fewer reward labels than baseline, same $\dgraph$) is a natural extra
axis for B3 that costs nothing extra to measure --- not yet run.
\end{itemize}

\subsection{B4: is the potential doing the work, or just any potential?}
\label{sec:b4-results}

\textbf{Methodology.}

\begin{algorithm}[H]
\caption{B4 methodology}
\begin{algorithmic}[1]
\State \textbf{Fixed:} \Call{Mode2-AuxRegression}{} loop (\cref{alg:consumption}), target form $-(1-\gamma^{d(s,g)})/(1-\gamma)$, $\lambda_{\text{aux}}=0.3$
\For{$d(s,g) \in \{$euclidean: $\|\zv_s-\zv_g\|$,\ graph: $\dgraph$,\ oracle: true wall-respecting pixel distance$\}$}
    \State train and evaluate as in \Call{Mode2-AuxRegression}{}, all six tiers, 3 seeds each
\EndFor
\State \textbf{Note:} euclidean is B0's known-bad proxy; oracle is not deployable (needs privileged position info) --- a noise-free ceiling
\end{algorithmic}
\end{algorithm}

\textbf{Results.}

\begin{table}[h]
\centering
\small
\caption{B4: peak held-out-test Spearman by distance source fed to mode 2, mean over 3 seeds.}
\begin{tabular}{lcccc}
\toprule
tier & baseline & euclidean & \textbf{graph} & oracle (ceiling) \\
\midrule
expert\_10   & $0.303$ & $0.314$ & $0.316$ & $0.452$ \\
expert\_25   & $0.197$ & $0.267$ & $\mathbf{0.351}$ & $0.435$ \\
expert\_50   & $0.272$ & $0.273$ & $\mathbf{0.383}$ & $0.562$ \\
expert\_100  & $0.306$ & $0.317$ & $\mathbf{0.646}$ & $0.728$ \\
mixed        & $0.340$ & $0.381$ & $\mathbf{0.490}$ & $0.641$ \\
mixed\_large & $0.575$ & $0.443$ & $\mathbf{0.908}$ & $0.931$ \\
\bottomrule
\end{tabular}
\end{table}

Graph beats euclidean at every tier, gap growing with data (mixed\_large: graph $0.908$, more
than double euclidean's $0.443$). Euclidean-mode-2 loses to baseline at mixed\_large ($0.443$
vs.\ $0.575$). Graph closes in on the oracle ceiling as data scales: the gap is largest at
expert\_50 ($0.383$ vs.\ $0.562$, a gap of $0.18$) and shrinks to $0.023$ at mixed\_large
($0.908$ vs.\ $0.931$).

\textbf{Takeaways.}
\begin{itemize}[leftmargin=1.4em,itemsep=2pt]
\item Graph clearly beats Euclidean everywhere, and the gap grows with data --- the core thesis
holds under the mechanism that actually works.
\item The auxiliary-regression mechanism is not sufficient by itself: regressing toward a bad
distance proxy (euclidean) can lose to doing nothing at all. It works specifically because the
graph gives it something accurate to regress toward.
\item Graph calibration quality itself improves with more data, converging toward the
ground-truth ceiling rather than plateauing below it.
\item Kill criterion (``raw Euclidean shaping matches graph shaping'') did not fire.
\end{itemize}

\subsection{B4 extension: transitions-only and edge-capped graphs}
\label{sec:b4ext-results}

\textbf{Methodology.}

\begin{algorithm}[H]
\caption{B4-extension methodology}
\begin{algorithmic}[1]
\State \textbf{Motivation:} a local, static (no RL) diagnostic campaign found that (i) same-episode ranking is measurably worse without identification edges, growing with data heterogeneity, and (ii) capping identification edges to each node's top-$k$ FAISS-HNSW neighbours costs nothing in Spearman quality (often slightly better) even at $k{=}4$ --- the fix B1 flagged as required before attempting Push-T (\cref{sec:pusht}). This extension checks both findings at the real RL-training level, not just the static metric.
\State \textbf{Fixed:} \Call{Mode2-AuxRegression}{} loop (\cref{alg:consumption}), identical target form and $\lambda_{\text{aux}}=0.3$, unchanged from B4
\For{$d(s,g) \in \{$transonly: $\dgraph$ with identification edges dropped,\ k4: $\dgraph$ capped to top-4 neighbours,\ k8: capped to top-8$\}$}
    \State train and evaluate as in \Call{Mode2-AuxRegression}{}, 4 tiers (expert\_50, expert\_100, mixed, mixed\_large), 5 seeds each
\EndFor
\State \textbf{Compute:} IIIT-H Ada cluster, direct-SSH orchestration (no SLURM), 60 runs total
\end{algorithmic}
\end{algorithm}

\textbf{Results.}

\begin{table}[h]
\centering
\small
\caption{B4-extension: peak held-out-test Spearman, mean $\pm$ std over 5 seeds (baseline and
full graph repeated from \cref{sec:b4-results} for reference, 3 seeds each).}
\begin{tabular}{lccccc}
\toprule
tier & baseline & graph (full) & transonly & \textbf{k4} & \textbf{k8} \\
\midrule
expert\_50   & $0.272$ & $0.383$ & $0.417 \pm 0.034$ & $\mathbf{0.500 \pm 0.040}$ & $0.463 \pm 0.037$ \\
expert\_100  & $0.306$ & $0.646$ & $0.646 \pm 0.014$ & $\mathbf{0.692 \pm 0.016}$ & $0.679 \pm 0.022$ \\
mixed        & $0.340$ & $0.490$ & $0.446 \pm 0.015$ & $\mathbf{0.584 \pm 0.023}$ & $0.575 \pm 0.023$ \\
mixed\_large & $0.575$ & $0.908$ & $0.790 \pm 0.015$ & $0.913 \pm 0.009$ & $\mathbf{0.915 \pm 0.009}$ \\
\bottomrule
\end{tabular}
\end{table}

\textbf{k4 and k8 do not merely match the full (uncapped) graph --- they beat it, at every tier
tested}: $+0.117$ / $+0.080$ at expert\_50, $+0.046$ / $+0.033$ at expert\_100, $+0.094$ / $+0.085$
at mixed, $+0.005$ / $+0.007$ at mixed\_large (k4/k8 respectively, vs.\ the full graph). This
directly extends the local static finding (capping costs nothing, often helps) to the metric that
actually matters.

\textbf{Transitions-only is competitive at the two pure-expert tiers but falls behind on the
mixed-quality ones}: statistically indistinguishable from the full graph at expert\_50
($+0.034$) and expert\_100 ($+0.000$), but clearly behind at mixed ($-0.044$) and mixed\_large
($-0.118$), the gap growing with data heterogeneity. This matches the local static
mechanism-probe's own pattern exactly (same-episode Spearman gap without identification edges was
smallest on pure-expert tiers, $+0.05$, and largest on mixed/mixed\_large, $+0.12$/$+0.14$) --- the
two independent diagnostics, one static and one from real training, agree on \emph{where}
identification-edge stitching matters most.

\textbf{Takeaways.}
\begin{itemize}[leftmargin=1.4em,itemsep=2pt]
\item Edge-count capping is not a quality trade-off to accept for Push-T-scale feasibility --- it
is a net improvement, confirmed now at the RL-training level, not just the static Spearman
diagnostic. This resolves B1's ``top-\,$k$ cap required before Push-T'' prerequisite with a
favourable answer rather than an open cost.
\item Cross-episode stitching's value is concentrated on heterogeneous data: on clean expert-only
tiers, a transitions-only graph is nearly as good as the full graph; on mixed-quality tiers, the
full graph's identification edges are clearly load-bearing, and the size of that gap tracks the
same pattern found statically outside any RL loop.
\item Practical implication for Push-T (\cref{sec:pusht}): a $k$-capped graph is not just feasible
at that scale, it is now the recommended default even where the uncapped graph would be
affordable, since it wins on Two-Room's own comparison.
\end{itemize}

\subsection{Live-rollout success rate: confirming the B3/B4 result}
\label{sec:actor-rollout}

\textbf{Methodology.}

\begin{algorithm}[H]
\caption{Live-rollout methodology}
\begin{algorithmic}[1]
\State \textbf{Actor loss (AWR):} $w \gets \mathrm{clip}\bigl(\exp(\beta(Q(s,a,g)-V(s,g))),\,100\bigr)$, $\beta=3.0$; $\ell_{\text{actor}} \gets w\cdot\|\pi(s,g)-a\|^2$; advantage computed under \texttt{no\_grad}
\State \textbf{peak\_rho:} cheap Spearman-based peak, tracked every 50 steps
\State \textbf{peak\_success:} live-rollout success rate over a fixed 12-episode subset, tracked every 2{,}500 steps, using \texttt{TwoRoomEnv}'s own criterion (distance-to-goal $<16$)
\State \textbf{Sweep:} 6 tiers $\times$ \{baseline, mode 2\} $\times$ 5 seeds $= 60$ runs, 20{,}000 steps each \Comment{truncated from 50{,}000 --- all pilot peaks landed by step 12{,}650}
\State \textbf{Final eval:} 100 episodes, at both the rho-selected and the success-selected checkpoint
\State \textbf{Compute:} IIIT-H Ada cluster
\end{algorithmic}
\end{algorithm}

\textbf{Results.}

\begin{table}[h]
\centering
\small
\caption{Live-rollout success rate, mean over 5 seeds, by checkpoint-selection criterion.}
\begin{tabular}{lcccc}
\toprule
tier & baseline (rho-sel) & baseline (success-sel) & \textbf{auxphi (rho-sel)} & \textbf{auxphi (success-sel)} \\
\midrule
expert\_10   & $0.414$ & $0.426$ & $0.406$ & $0.420$ \\
expert\_25   & $0.480$ & $0.458$ & $0.490$ & $\mathbf{0.514}$ \\
expert\_50   & $0.464$ & $0.406$ & $\mathbf{0.476}$ & $\mathbf{0.508}$ \\
expert\_100  & $0.588$ & $0.620$ & $\mathbf{0.712}$ & $\mathbf{0.714}$ \\
mixed        & $0.402$ & $0.422$ & $0.448$ & $\mathbf{0.472}$ \\
mixed\_large & $0.668$ & $0.718$ & $\mathbf{0.788}$ & $\mathbf{0.808}$ \\
\bottomrule
\end{tabular}
\end{table}

Auxiliary regression beats baseline at 6 of 6 tiers under both selection criteria (up from 5 of 6
on the proxy metric --- expert\_10 was a genuine wash there, a small consistent win here). Margin
grows with data exactly as the Spearman result predicted: negligible at expert\_10, $+0.09$ to
$+0.19$ by expert\_100 and mixed\_large. Selecting by the success peak gives a small edge over
selecting by the rho peak in 10 of 12 rows; \texttt{peak\_success} is often first observed at
step 2{,}500 while \texttt{peak\_rho} can be as early as step 300.

\textbf{Takeaways.}
\begin{itemize}[leftmargin=1.4em,itemsep=2pt]
\item The proxy-metric result converts into real task success: this is direct evidence that the
value-ranking win measured throughout B3--B4 was tracking something that mattered for the policy
actually extracted from it, not an artifact of the ranking metric.
\item The two peaks (rho-based, success-based) genuinely track different things during training
and land at different steps --- justifying tracking both rather than assuming they coincide.
\end{itemize}

\subsection{Live-rollout success rate: confirming the B4 extension}
\label{sec:b4ext-actor-rollout}

\textbf{Methodology.}

\begin{algorithm}[H]
\caption{B4-extension live-rollout methodology}
\begin{algorithmic}[1]
\State \textbf{Fixed:} identical actor (AWR) and dual peak-tracking setup as \cref{sec:actor-rollout} --- only the distance source fed to mode 2 changes
\For{$d(s,g) \in \{$transonly, k4, k8$\}$ (as defined in \cref{sec:b4ext-results})}
    \State train actor, live-rollout eval as in \cref{sec:actor-rollout}, 4 tiers (expert\_50, expert\_100, mixed, mixed\_large), 5 seeds each
\EndFor
\State \textbf{Compute:} IIIT-H Ada cluster, direct-SSH orchestration, 60 runs (train $\to$ rho-eval $\to$ success-eval chained per run)
\end{algorithmic}
\end{algorithm}

\textbf{Results.}

\begin{table}[h]
\centering
\small
\caption{Live-rollout success rate at the success-selected checkpoint, mean $\pm$ std over 5
seeds (baseline and full graph repeated from \cref{sec:actor-rollout} for reference).}
\begin{tabular}{lccccc}
\toprule
tier & baseline & graph (full) & transonly & \textbf{k4} & \textbf{k8} \\
\midrule
expert\_50   & $0.406$ & $0.508$ & $0.410 \pm 0.025$ & $0.448 \pm 0.061$ & $\mathbf{0.492 \pm 0.033}$ \\
expert\_100  & $0.620$ & $0.714$ & $0.538 \pm 0.038$ & $\mathbf{0.694 \pm 0.023}$ & $0.692 \pm 0.015$ \\
mixed        & $0.422$ & $0.472$ & $0.380 \pm 0.044$ & $\mathbf{0.472 \pm 0.065}$ & $0.460 \pm 0.046$ \\
mixed\_large & $0.718$ & $0.808$ & $0.686 \pm 0.038$ & $\mathbf{0.804 \pm 0.020}$ & $0.792 \pm 0.046$ \\
\bottomrule
\end{tabular}
\end{table}

Unlike the Spearman-proxy result (\cref{sec:b4ext-results}), where k4 and k8 beat the full
uncapped graph outright at every tier, on real rollout success they instead \emph{match} it
closely: k4 ties the full graph exactly at mixed ($0.472$ vs.\ $0.472$) and comes within $0.02$ of
it at expert\_100 and mixed\_large, only clearly trailing at expert\_50 ($0.448$ vs.\ $0.508$).
Both k4 and k8 still beat baseline decisively at all four tiers ($+0.04$ to $+0.09$).
Transitions-only tells the more consequential story: the proxy metric called it ``competitive at
the two pure-expert tiers,'' but on real rollout success it is at or below baseline at all four
(a wash at expert\_50, $-0.04$ to $-0.08$ elsewhere) --- it never wins. Selecting the checkpoint by
the rho-peak instead of the success-peak gives the same qualitative picture (k4/k8 close to the
full graph and clearly ahead of baseline; transonly at or below baseline everywhere), so this is
not an artifact of the checkpoint-selection criterion.

\textbf{Takeaways.}
\begin{itemize}[leftmargin=1.4em,itemsep=2pt]
\item The recommendation to use a $k$-capped graph as the default construction method
(\cref{sec:b4ext-results}) survives contact with real task performance: k4/k8 remain a clear,
consistent win over baseline, even though their edge over the \emph{full} graph specifically was a
proxy-metric artifact that does not reproduce at the real-task level.
\item Transitions-only's failure mode is worse in practice than the proxy metric suggested: it
does not merely fall behind the full graph on heterogeneous data, it fails to beat doing nothing
(baseline) at any tier tested. Identification edges are not optional at the real-task level, even
on the pure-expert tiers where the static and Spearman-proxy diagnostics suggested near-parity.
\item General lesson for this project: the Spearman-vs-oracle proxy metric is directional but not
quantitatively reliable close to the top of the ranking --- differences among the strongest few
conditions (full graph vs.\ k4 vs.\ k8) do not survive to real task success in the same order or
magnitude, even though the much larger baseline/transonly-vs-graph gaps do.
\end{itemize}

\section{Push-T: transferring the approach to a second environment}
\label{sec:pusht}

Everything above is Two-Room. This section repeats the same questions --- does the calibrated
graph beat raw latent distance, does it help offline RL, does it hold up under a live-rollout
success check --- on Push-T, a contact-manipulation task with a very different observation and
dynamics structure. It is reported as its own arc rather than folded into the Two-Room ladder
above, because it did not transfer cleanly on the first attempt: the calibration formula fails
outright, the live-rollout protocol initially used was too hard to show any signal at all, and
diagnosing the resulting apparent underperformance took several rounds of investigation. Each
subsection below is a self-contained result with its own methodology and configuration; later
subsections build on fixes established in earlier ones, noted explicitly where relevant.

\subsection{Does the identification-edge calibration transfer?}
\label{sec:b5-results}

\textbf{Methodology.}

\begin{algorithm}[H]
\caption{Push-T gate-check methodology --- Two-Room's B0 gate, ported to Push-T}
\begin{algorithmic}[1]
\State \textbf{Checkpoint:} real pretrained \texttt{quentinll/lewm-pusht} (downloaded fresh, not the local from-scratch epoch-1 checkpoint, which is confirmed collapsed --- \cref{sec:b2-results}-style off-diagonal cosine and participation-ratio check run first)
\State \textbf{Landmarks:} 300 episodes ($36{,}824$ frames), encoded with the same frozen-encoder pipeline as \cref{alg:construction}
\State \textbf{Oracle:} no walls to route around, so ground truth is Euclidean distance in z-scored privileged state $[\text{agent}_{x,y}, \text{block}_{x,y}, \cos\theta_{\text{block}}, \sin\theta_{\text{block}}]$, not a navigation-grid distance
\State \textbf{Calibration:} hold $q=10^{-3}$ fixed at the Two-Room value, as proposed
\If{calibrated $\varepsilon^2 \le 0$ (rho\_hat statistically $\ge 1$)}
    \State \textbf{fall back} to the named contingency: heuristic (uncalibrated) top-$k$ identification edges, $k=4$ carried over from \cref{sec:b4ext-results} unchanged
\EndIf
\State compare Spearman(oracle, graph-geodesic) vs.\ Spearman(oracle, raw latent Euclidean)
\end{algorithmic}
\end{algorithm}

\textbf{Results.}

The calibrated threshold failed outright, not marginally: $\hat\rho$ came out statistically
$\ge 1$ (measured $1.0009$ at 300 episodes, up from an already-pathological $0.9985$ at 20 ---
more landmarks made it \emph{worse}, ruling out a landmark-coverage explanation), forcing
$\varepsilon^2 < 0$ and exactly zero identification edges at every quantile tested, at both
landmark scales. This is not a checkpoint-collapse artifact: a properly-randomized,
large-sample participation-ratio check (the naive small-batch check is biased by temporally
correlated sampling and understates it) puts this checkpoint's effective rank at $84$ of $192$
dims, at least as healthy as Two-Room's $37$ of $192$. The real cause is narrower: adjacent-frame
latent displacement is tiny relative to the cross-episode spread on Push-T (squared-distance
ratio $\approx 176\times$) versus Two-Room's $\approx 2.4\times$ --- Two-Room's calibration
formula implicitly assumes a margin this checkpoint's geometry does not have, independent of its
overall quality.

\begin{table}[h]
\centering
\small
\caption{Push-T: Spearman vs.\ the privileged-state oracle, 300-episode landmark set,
heuristic top-4 identification edges (calibrated threshold failed --- see above).}
\begin{tabular}{lcc}
\toprule
identification-edge weight & raw euclidean & \textbf{graph geodesic} \\
\midrule
$0.0$ & $0.219$ & $0.319$ \\
$0.5$ & $0.219$ & $0.518$ \\
$1.0$ & $0.219$ & $0.537$ \\
$2.0$ & $0.219$ & $\mathbf{0.540}$ \\
$4.0$ & $0.219$ & $0.534$ \\
\bottomrule
\end{tabular}
\end{table}

With the heuristic fallback, the graph beats raw Euclidean by more than $2\times$ ($0.540$ vs.\
$0.219$ at the best identification-edge weight) --- a clear pass of the same gate B0 applied to
Two-Room, on $98.3\%$ of sampled pairs (up from under $2\%$ finite under the failed calibrated
graph, whose near-total disconnection made most cross-episode comparisons undefined).

\textbf{Takeaways.}
\begin{itemize}[leftmargin=1.4em,itemsep=2pt]
\item The kill criterion (``needs per-environment retuning'') fired in its precise, named
form --- the calibration quantile did not transfer as-is --- but the documented fallback (a
heuristic, uncalibrated top-$k$ graph) recovers a decisive win over raw Euclidean anyway. This is
the weaker-but-still-useful deployment story the proposal anticipated, not a failure of the core
idea.
\item The failure mode is specific and diagnosable, not a vague ``it didn't work'': it is a
mismatch between how large adjacent-frame latent movement is relative to the overall latent
spread on the two environments, not a difference in checkpoint quality (which is comparable or
better on Push-T by the properly-measured participation ratio).
\item $k=4$ was carried over unchanged from Two-Room's B4-extension winner, not retuned for
Push-T --- a natural next check, not yet done.
\item This result used only the local machine's GPU (300-episode encode $+$ graph $+$ eval, a
few minutes total) --- no Ada allocation needed for B0. The RL-training pipeline result below
did need Ada.
\end{itemize}

\subsection{Does the graph help offline RL training?}
\label{sec:b5-b3-results}

\textbf{Methodology.}

\begin{algorithm}[H]
\caption{Push-T RL-training methodology --- Two-Room's B3 data-tier sweep, ported to Push-T}
\begin{algorithmic}[1]
\State \textbf{Fixed:} \Call{Mode2-AuxRegression}{} loop (\cref{alg:consumption}), identical target form and $\lambda_{\text{aux}}=0.3$, unchanged from Two-Room's B3/B4 --- \texttt{run\_condition\_resumable} reused with zero code changes, only a Push-T-specific \texttt{setup} dict
\State \textbf{Graph:} heuristic top-4 identification edges (\cref{sec:b5-results}), built fresh per tier
\State \textbf{Tiers:} expert\_10/25/50/100 (real expert episodes only) and mixed (50 real expert $+$ 50 WeakPolicy-rollout episodes --- stable\_worldmodel's built-in weak/unfocused Push-T collection policy, validated to fail the task: 0/20 rollouts reached the goal within 250 steps at every tested setting, confirming it is genuinely low-quality data, not disguised expert data)
\State train and evaluate baseline ($\lambda_{\text{aux}}=0$) and auxphi ($\lambda_{\text{aux}}=0.3$), all 5 tiers, 3 seeds each, 50{,}000 steps
\State \textbf{Compute:} IIIT-H Ada cluster, direct-SSH orchestration, 30 runs total
\end{algorithmic}
\end{algorithm}

\textbf{Results.}

\begin{table}[h]
\centering
\small
\caption{Push-T: peak held-out-test Spearman, mean over 3 seeds.}
\begin{tabular}{lccc}
\toprule
tier & baseline & \textbf{auxphi} & $\Delta$ \\
\midrule
expert\_10  & $0.148$ & $\mathbf{0.178}$ & $+0.030$ \\
expert\_25  & $0.284$ & $\mathbf{0.330}$ & $+0.046$ \\
expert\_50  & $0.128$ & $\mathbf{0.182}$ & $+0.054$ \\
expert\_100 & $0.074$ & $\mathbf{0.144}$ & $+0.070$ \\
mixed       & $0.391$ & $\mathbf{0.418}$ & $+0.027$ \\
\bottomrule
\end{tabular}
\end{table}

Auxphi beats baseline at all 5 of 5 tiers, matching Two-Room's B3/B4 finding qualitatively. Two
things do \emph{not} match Two-Room, and are reported rather than smoothed over: (1) absolute
Spearman is far lower here (peaking at $0.42$ vs.\ Two-Room's $0.91$ at its best tier) and (2) it
is \emph{non-monotonic} in tier size --- expert\_25 ($0.284$/$0.330$) outscores both expert\_50
and expert\_100, the opposite of Two-Room's clean monotonic improvement with data. This is not a
graph-construction artifact: identification-edge count and average degree both scale up sensibly
with landmark count (degree $0.05\to0.07\to0.34\to0.50\to0.44$ across
expert\_10/25/50/100/mixed), the same qualitative pattern B1 found healthy on Two-Room. The more
likely cause is a fixed training budget (50{,}000 steps, unchanged from Two-Room) not scaling
with the HER-tuple pool size, which grows faster on Push-T than on Two-Room at the same nominal
episode count (Push-T's episodes average $125$ steps vs.\ Two-Room's $92$) --- larger Push-T
tiers may simply be undertrained relative to smaller ones at a fixed step count, and/or per-tier
oracle z-scoring (recomputed fresh per tier, unlike Two-Room's fixed wall-geometry oracle) makes
absolute values not strictly comparable across tiers. Neither is confirmed; both are candidates
for a longer training-budget or fixed-oracle-normalization follow-up rather than settled fact.

\textbf{Takeaways.}
\begin{itemize}[leftmargin=1.4em,itemsep=2pt]
\item The core claim replicates on a second environment: the graph-based auxiliary-regression
mechanism beats baseline everywhere tested, using the same code, the same mechanism, and (for
identification edges) the same fallback construction method established in \cref{sec:b5-results}.
\item Absolute magnitude and its scaling with data size do \emph{not} replicate cleanly --- this
is Push-T's first RL-training result, not yet debugged to the same depth as Two-Room's multi-round
B3/B4 process was, and the non-monotonic tier-size pattern is flagged as open rather than
explained away.
\item Live-rollout confirmation (the real task-success check that mattered for Two-Room's
headline result, \cref{sec:actor-rollout}) has not yet been built for Push-T --- this proxy-metric
result alone should be read with the same caution the B4-extension proxy result eventually
needed (\cref{sec:b4ext-actor-rollout}).
\end{itemize}

\subsection{Diagnosing why the calibrated threshold fails}
\label{sec:b5-calib-diagnosis}

\textbf{Root cause.} The calibration formula $\hat\rho = \tfrac{1}{d}\mathbb E[\zv_t^\top
\zv_{t+1}]$ (mean inner product between temporally-adjacent same-episode frames, normalized by
the raw embedding dimension $d{=}192$) implicitly assumes two things: that $\mathbb
E[\|\zv\|^2] = d$, and that the embedding's effective degrees of freedom equal $d$. Neither
holds tightly enough on Push-T. Measured directly: $\mathbb E[\|\zv\|^2] = 193.2$ (a $0.6\%$
deviation from the assumed $192$), and the covariance's participation ratio --- effective
dimensionality, not the raw ambient one --- is $84$ of $192$ dims. On Two-Room ($\hat\rho{=}0.584$)
this $0.6\%$ normalization error is negligible relative to $1-\hat\rho \approx 0.42$; on Push-T
($1-\hat\rho \approx 0.001$) the same error is large enough to flip the sign, since $\varepsilon^2
= 2(1-\hat\rho)F^{-1}_{\chi^2_d}(q)$ collapses to a negative ``variance'' the moment $\hat\rho \ge
1$. This is not a sign the checkpoint is worse on Push-T: an independently-measured, large-sample
participation ratio puts it at $84$ of $192$ effective dims, comparable to Two-Room's own. The
formula's assumptions break, not the representation.

\textbf{Is a different distance metric the fix?} Tested directly rather than assumed: dot
product, cosine similarity (raw and mean-centered), per-dimension standardization, and a
PCA-whitened basis, each scored as edge precision (fraction of $k{=}4$ cross-episode nearest
neighbors whose true state distance is under the same-state threshold $\tau$) against the raw
Euclidean baseline already in use.

\begin{table}[h]
\centering
\small
\caption{Candidate identification-edge metrics, scored against the ground-truth state oracle
(300-episode Push-T landmark set, $n{=}36{,}824$).}
\begin{tabular}{lcc}
\toprule
metric & Spearman(true, metric) & edge precision @$k{=}4$ \\
\midrule
raw Euclidean (current) & $0.174$ & $0.576$ \\
PCA-whitened & $0.204$ & $0.579$ \\
per-dim standardized & $0.147$ & --- \\
centered cosine & $0.050$ & --- \\
cosine / dot product & $0.037$ & --- \\
\bottomrule
\end{tabular}
\end{table}

Cosine/dot-product similarity is \emph{worse} than raw Euclidean, not better; whitening is a
wash ($0.579$ vs.\ $0.576$ precision, well within noise). The metric was never the problem ---
Euclidean distance in this latent already separates near/far states about as well as any linear
transform of it can.

\textbf{Precision vs.\ connectivity: the real tradeoff.} A corrected calibration formula ---
$\varepsilon^2 = \overline{\|\zv_t - \zv_{t+1}\|^2}\,/\,d_{\text{eff}} \cdot
F^{-1}_{\chi^2_{d_{\text{eff}}}}(q)$, using the measured adjacent-frame squared displacement
directly and the participation ratio $d_{\text{eff}}$ in place of $d$ --- is always positive by
construction (both factors are non-negative), giving $\varepsilon^2{=}1.32$ at $q{=}10^{-3}$.
Applying it raises identification-edge precision from $57.6\%$ to $96.6\%$. But the B0 gate
metric gets \emph{worse}, not better:

\begin{table}[h]
\centering
\small
\caption{B0 gate Spearman under different identification-edge thresholds, same construction
otherwise (300-episode landmark set, $k{=}4$-capped kNN, identification-edge weight $1.0$).}
\begin{tabular}{lcccc}
\toprule
threshold & $\varepsilon^2$ & id.\ edges & unreachable eval pairs & Spearman(true, graph) \\
\midrule
$\varepsilon^2=\infty$ (prior fallback) & $\infty$ & $72{,}722$ & $17.3\%$ & $\mathbf{0.553}$ \\
corrected, $q{=}10^{-3}$ & $1.32$ & $16{,}072$ & $99.6\%$ & $0.058$ \\
adjacent-pair empirical $q_{0.9}$ & $5.21$ & $32{,}463$ & $98.0\%$ & $0.057$ \\
adjacent-pair empirical $q_{0.99}$ & $14.16$ & $49{,}297$ & $77.5\%$ & $0.191$ \\
adjacent-pair empirical $q_{0.999}$ & $27.20$ & $62{,}333$ & $40.5\%$ & $0.396$ \\
\bottomrule
\end{tabular}
\end{table}

Spearman tracks connectivity, not precision: the correctly-calibrated, high-precision threshold
shatters the graph into fragments and the gate metric collapses below the no-graph baseline
($0.217$ raw Euclidean). Precision-weighted edge cost (making a distant edge expensive rather
than dropping it, keeping every $k{=}4$ candidate) was also tested across eight weighting
schemes and never beat flat weighting on the already-dense ($\varepsilon^2{=}\infty$) graph ---
confirming connectivity, not edge confidence, is the dominant factor once a graph is already
connected.

\textbf{Why the graph fragments.} At $q{=}10^{-3}$ the corrected threshold produces $91$
connected components from $300$ episodes: $34$ episodes end up fully isolated (no cross-episode
identification edge at all), $83$ episodes merge into one giant component, and the remaining
$183$ scatter across $56$ smaller clusters. This is not primarily a data-sparsity problem ---
$68.5\%$ of all landmarks genuinely have a true same-state twin somewhere in the other $299$
episodes --- it is a \emph{recall} problem specific to what this latent encodes. Of landmarks
with a genuine true-state twin, the strict threshold finds only $38.7\%$; widening the kNN
search from $k{=}4$ to $k{=}1024$ recovers none of the rest (edge count and Spearman both
saturate by $k{\approx}64$), meaning the missed matches are not merely ranked outside the top
few neighbors, they are placed genuinely far away in latent space. The missing $61.3\%$
correlates cleanly with a variable the oracle never sees: agent/block \emph{velocity} (state
columns excluded from the position-and-angle oracle). Splitting by velocity-gap quartile
between a landmark and its true partner: $P(\text{match found}) = 64.9\%, 47.0\%, 27.0\%,
15.9\%$ (Spearman $=-0.40$, $p\approx0$). LeWM's latent is trained on a next-frame predictive
objective, so two frames sharing position and angle but differing in momentum are, correctly
from the encoder's own perspective, headed toward different futures --- the graph is answering a
stricter, more predictively-faithful question (``same state, including momentum'') than the one
the downstream task actually asks (``same useful position''; the paper's own live-rollout
protocol and this oracle both define goals by position and angle alone, never velocity).

\textbf{Does a fragmented graph degrade the RL mechanism to imitation?} No, for a structural
reason, not a lucky empirical one. HER goals are always drawn from the \emph{same} episode's own
future (\texttt{build\_her\_tuples}); the base $-1/0$ TD objective never reads the graph, and the
auxiliary regression target queries $\Phi(s,g)$ only on those same-episode pairs, which a chain
of directed transition edges always connects regardless of whether any identification edge
exists at all --- cross-episode fragmentation is structurally invisible to what actually drives a
training step. The real risk is narrower: for an episode with no useful cross-episode shortcut,
$\Phi(s,g)$ degrades toward that episode's own chronological step count, close to redundant with
what the $-1/0$ reward already encodes, so the auxiliary loss's marginal contribution shrinks
toward the baseline rather than the mechanism collapsing outright.

\textbf{Comparing the stitching benefit to Two-Room's.} Restricting Two-Room's own measured
same-episode Spearman gain from identification edges (\cref{sec:b1-results}) to a like-for-like
comparison: Push-T's gain is real but structurally smaller, and for a specific, checkable reason
--- Two-Room's graph is always one connected component drawing shortcuts from all $300$ episodes,
while even Push-T's best-connected episodes only draw from whichever cluster ($2$ to $83$
episodes) they happen to land in.

\begin{table}[h]
\centering
\small
\caption{Same-episode Spearman gain from identification edges, Push-T by component membership.
Two-Room, for comparison, is always a single connected component and measures $+0.046$ to
$+0.143$ across its 6 tiers (\cref{sec:b1-results}).}
\begin{tabular}{lc}
\toprule
group & Push-T gain \\
\midrule
isolated episodes ($34$ of $300$) & $+0.015$ \\
giant component ($83$ episodes) & $+0.036$ \\
small merged clusters ($183$ episodes across $56$ components) & $+0.066$ \\
\bottomrule
\end{tabular}
\end{table}

Isolated and giant-component episodes both fall below Two-Room's weakest tier; the small merged
clusters land inside Two-Room's range, in some cases above its median tier. The mechanism is the
same one validated on Two-Room --- real cross-episode shortcut-stitching, not merely
connectivity-patching --- it is just capped by a smaller shortcut pool on Push-T, on top of (not
instead of) the velocity-driven precision/recall tradeoff above.

\textbf{Takeaways.}
\begin{itemize}[leftmargin=1.4em,itemsep=2pt]
\item The calibration failure has a specific, checkable cause (a normalization assumption
violated by $\approx0.6\%$ that only matters because $1-\hat\rho$ is itself tiny on Push-T), not
a vague ``it doesn't transfer.''
\item No alternative distance metric (dot product, cosine, whitening) closes the gap --- ruled
out directly, not assumed.
\item Precision and connectivity trade off sharply on this environment specifically: the
statistically correct threshold is a \emph{worse} graph by the gate's own metric, because it
shatters connectivity faster than it improves confidence.
\item The recall gap driving fragmentation is attributable, in large part, to velocity --- a
variable this predictive latent encodes but the position-only oracle does not ask about. Whether
to fix the oracle (add velocity) or accept this as inherent to a predictive latent is an open
design choice, not yet resolved.
\item The RL-training mechanism (\cref{sec:method}) is structurally immune to cross-episode
fragmentation by construction (HER is always same-episode); what degrades under fragmentation is
the auxiliary signal's marginal value on affected episodes, not training correctness.
\end{itemize}

\subsection{A corrected calibration formula, and a B0/B1/B3-style re-run}
\label{sec:b5-calib-fix}

\textbf{The fix.} \texttt{compute\_calibration()} (\texttt{scripts/common/graph\_lib.py}) now
computes $\varepsilon^2$ directly from the measured adjacent-pair squared displacement and the
covariance's participation ratio, replacing the $\hat\rho$/$d$-based formula
(\cref{sec:b5-calib-diagnosis}). This is always positive by construction, so the
$\varepsilon^2{=}\infty$ fallback --- which was never really an approximation of a thresholded
search, since it accepts a node's $k$ nearest neighbors however far away they are --- is removed
entirely rather than special-cased away. \texttt{build\_id\_edges\_faiss\_capped}'s kNN search is
unchanged: it was already correct (every edge it returns also satisfies the threshold, so its
output was always a subset of the exact brute-force scan's), the fallback was the only path that
broke that guarantee. The fix is scoped to \texttt{compute\_calibration()} specifically, which
only the ``capped'' construction path (currently Push-T only) calls; Two-Room's ``bruteforce''
path (\texttt{calibrate()}) is left untouched to preserve its validated bit-exact historical
numbers --- confirmed unchanged ($\varepsilon^2{=}114.10$) after the change. One narrower side
effect: Two-Room's B4-extension $k\{K\}$ distance-source ablation arms also call
\texttt{compute\_calibration()}, so $\varepsilon^2$ shifts there too ($114.10 \to 94.06$); the
main graph condition and headline B4/B4-extension results (\cref{sec:b4ext-results}) are
unaffected.

\textbf{B0, re-run.} Confirms the diagnosis on the live production path rather than a scratch
script: $\hat\rho{=}1.0009$, $\varepsilon^2{=}1.3226$, $13{,}750$ identification edges, and the
gate's own evaluation collapses to $16$ of $4{,}000$ sampled pairs finite (\texttt{outputs/
b0\_pusht/b0\_results.json}). The fix delivers a principled, always-well-defined,
subset-preserving construction exactly as specified --- it does not, by itself, resolve the
connectivity problem diagnosed above.

\textbf{B1, written and run for Push-T.} No unified Push-T B1 existed before this; adapted from
Two-Room's script onto the corrected \texttt{common.graph\_lib} layer, testing the $k$-capped
construction actually deployed (Two-Room's B1 tested exact brute force, which is the thing B1
itself found infeasible at scale). Four findings:
\begin{itemize}[leftmargin=1.4em,itemsep=2pt]
\item \textbf{Edge count scales as $n^{1.61}$} (50/300/1000-episode landmark sets), far better
than Two-Room's uncapped $n^{1.96}$ but not linear either --- average degree itself grows with
scale ($0.26 \to 0.75 \to 1.56$), since a larger dataset increases the odds any given landmark
has a genuine revisit.
\item \textbf{The successor-consistency filter (\cref{sec:b1-results}), which did not help on
Two-Room, looks promising here}: of $13{,}751$ identification edges, $82.5\%$ were testable and
$86.1\%$ passed; dropping the $1{,}583$ that failed took the (still small-sample) gate Spearman
from $0.897$ to $0.975$. Not yet validated at a sample size large enough to trust the magnitude.
\item \textbf{Fragmentation is graded across the full $q$ range, not a cliff at $q{=}10^{-3}$}:
even the loosest $q{=}0.5$ tested gives only $44.7\%$ largest-component coverage and $24.4\%$
cross-episode reachability; both fall monotonically to single digits by $q{=}10^{-9}$.
\item \textbf{No pathological hub nodes} (max degree $10$, structurally bounded near $2k$) ---
the $k$-cap does what it is meant to.
\end{itemize}

\textbf{B3, re-run.} The prior $30$-run sweep (5 tiers $\times$ 2 conditions $\times$ 3 seeds,
\cref{sec:b5-b3-results}) was built entirely under the removed $\varepsilon^2{=}\infty$
fallback; those checkpoints and the stale per-tier $\Phi$ cache were preserved
(\texttt{outputs/b3\_pusht/\{checkpoints,tier\_cache\}\_pre\_calib\_fix/}) rather than
overwritten, and the sweep re-launched from scratch under the corrected, fragmented-but-principled
graph, same tiers/seeds/step budget.

\begin{table}[h]
\centering
\small
\caption{B3 re-run under the corrected calibration: peak held-out-test Spearman, mean over 3
seeds. Baseline never reads the graph, so it serves as a same-seed control between the two
sweeps.}
\begin{tabular}{lcccc}
\toprule
tier & baseline (prior $\to$ re-run) & auxphi (prior $\to$ re-run) & $\Delta$ (prior) & $\Delta$ (re-run) \\
\midrule
expert\_10  & $0.148 \to 0.1476$ & $0.178 \to \mathbf{0.1772}$ & $+0.030$ & $+0.030$ \\
expert\_25  & $0.284 \to 0.2838$ & $0.330 \to \mathbf{0.3255}$ & $+0.046$ & $+0.042$ \\
expert\_50  & $0.128 \to 0.1282$ & $0.182 \to \mathbf{0.1745}$ & $+0.054$ & $+0.046$ \\
expert\_100 & $0.074 \to 0.0738$ & $0.144 \to \mathbf{0.1293}$ & $+0.070$ & $+0.056$ \\
mixed       & $0.391 \to 0.3909$ & $0.418 \to \mathbf{0.4638}$ & $+0.027$ & $+0.073$ \\
\bottomrule
\end{tabular}
\end{table}

Baseline is essentially unchanged tier-for-tier ($\le0.001$ everywhere) --- exactly what the
mechanism argument above predicts, since baseline never reads the graph and the two sweeps
share seeds; this is a useful same-seed control confirming nothing else about the setup moved
between runs. Auxphi still beats baseline at all $5$ of $5$ tiers under the far more fragmented
graph, but the per-tier margin is not uniformly larger or smaller: it shrinks at
expert\_25/50/100 (most at expert\_100, $+0.070\to+0.056$) and grows sharply at mixed
($+0.027\to+0.073$, now easily the largest margin of any tier in either sweep). Connectivity
collapsing from $82.7\%$ reachable (\cref{sec:b5-calib-diagnosis}'s $\varepsilon^2{=}\infty$ row)
to $0.4\%$ reachable did not break the auxiliary signal, consistent with the structural argument
that same-episode HER pairs never depend on cross-episode connectivity --- but it did not leave
tier-by-tier magnitudes untouched either, and the shift has no settled explanation yet: mixed's
combination of real and noisy episodes may interact with the corrected graph's episode-loop-closure
edges (\cref{sec:b5-calib-diagnosis}) differently than expert-only tiers do, but this is a
hypothesis, not a confirmed mechanism.

\subsection{The live-rollout evaluation protocol}
\label{sec:b5-eval-protocol}

\textbf{Motivation.} Every Push-T actor described so far was evaluated on arbitrary
cross-episode $(s,g)$ pairs from the B0/B1 Spearman eval set, with a $250$-step budget
(\texttt{eval\_protocol=cross\_episode}) --- the original convention, ported unchanged from
Two-Room. Every condition scored $0.000$ under it (\cref{sec:b5-b3-results}). Before reading
that as ``the method fails on Push-T,'' the protocol itself needed checking against how the
LeWM paper actually measures its own Push-T policies (Fig.~6: GCBC~$75\%$, GCIVL~$33\%$,
GCIQL~$20\%$, Random~$2\%$): \texttt{stable-worldmodel}'s \texttt{scripts/plan/eval\_ff.py}
(App.~F.1) starts from a random dataset state, sets the goal to the state exactly $25$
timesteps later in the \emph{same} trajectory, and allows $50$ env steps to reach it ---
success is the environment's own \texttt{terminated} flag.

\begin{algorithm}[H]
\caption{\texttt{eval\_protocol=same\_episode} (\texttt{config/graph/actor.yaml}, now default)}
\begin{algorithmic}[1]
\State start $\gets$ a row sampled from \texttt{eval\_pool} (\texttt{dataset} = the full h5
minus the tier's train episodes, or \texttt{test\_episodes} = the tier's own held-out episodes)
\State goal $\gets$ the state \texttt{env.paper\_goal\_offset}$=25$ steps later in the same episode
\State run the actor for up to \texttt{env.paper\_eval\_budget}$=50$ steps; success = env's own \texttt{eval\_state()}
\end{algorithmic}
\end{algorithm}

\textbf{Results.} All runs below use the exact same actor-training loop and hyperparameters as
Two-Room's B3/B4 (\texttt{aux\_lambda=0.3}, \texttt{her\_goal\_gamma=0.9},
\texttt{aux\_standardize=false}, \texttt{awr\_normalize\_adv=false}, \texttt{aux\_cross\_lambda=0.0}
--- i.e.\ ``vanilla'' auxphi, none of the Push-T-specific knobs introduced in
\cref{sec:b5-underperf-diagnosis} below), trained via \texttt{scripts/actor\_train.py env=pusht
tier=<tier> variant=<baseline|auxphi> seed=0} and scored via
\texttt{scripts/actor\_rollout\_eval.py} with the config shown above, \texttt{eval\_seed=42},
$n=100$ episodes (random baseline: $n=20$, \texttt{policy=random}, no checkpoint).

\begin{table}[h]
\centering
\small
\caption{Live-rollout success rate under the paper's own protocol
(\texttt{eval\_protocol=same\_episode}, \texttt{eval\_pool=dataset}, \texttt{eval\_seed=42}).}
\begin{tabular}{llccc}
\toprule
tier & condition & select & $n$ & success rate \\
\midrule
expert\_25 (20 train eps) & random policy & --- & 20 & $15\%$ \\
expert\_25 & baseline & rho & 100 & $7\%$ \\
expert\_25 & \textbf{auxphi} & rho & 100 & $7\%$ \\
expert\_25 & baseline & success & 100 & $7\%$ \\
expert\_25 & \textbf{auxphi} & success & 100 & $7\%$ \\
expert\_25 & baseline, $\gamma_{\text{HER}}=0.96$ & rho & 100 & $2\%$ \\
expert\_25 & baseline, $\gamma_{\text{HER}}=0.96$ & success & 100 & $5\%$ \\
expert\_100 (80 train eps) & baseline & rho & 100 & $22\%$ \\
expert\_100 & \textbf{auxphi} & rho & 100 & $\mathbf{26\%}$ \\
\bottomrule
\end{tabular}
\end{table}

\textbf{Takeaways.}
\begin{itemize}[leftmargin=1.4em,itemsep=2pt]
\item The $0.000$ result was a protocol artifact, not a method failure: the same expert\_100
checkpoints that scored $0$ under \texttt{cross\_episode} score $22$--$26\%$ under the paper's
own protocol, comfortably above the measured random floor ($15\%$ on this pool; the next
subsection reports a lower $4$--$9\%$ floor on other pools --- the floor itself is
pool-dependent and must be measured alongside any headline number, not assumed).
\item At expert\_25, baseline and auxphi are statistically indistinguishable (both $7\%$, both
selection rules) --- essentially tied at the random floor with only $20$ train episodes.
Widening the HER goal horizon ($\gamma_{\text{HER}}: 0.9\to0.96$) did not help; it made the
baseline worse ($7\%\to2$--$5\%$), so the goal-sampling distribution is not the bottleneck at
this data scale.
\item At expert\_100, auxphi's edge over baseline ($26\%$ vs.\ $22\%$) reappears once there is
enough data to separate the two conditions from noise --- consistent with, though smaller in
magnitude than, the Spearman-proxy gap measured for the same tier
(\cref{sec:b5-b3-results}: $+0.070$).
\item Data scale, not any of the algorithmic knobs tried so far, is the dominant lever toward
the paper's own $20$--$75\%$ range --- see the next subsection.
\end{itemize}

\subsection{Diagnosing the apparent underperformance}
\label{sec:b5-underperf-diagnosis}

\textbf{Two plausible explanations, tested directly and refuted.} Diagnostics in
\texttt{scripts/investigations/pusht\_lowscore/\{diagnose.py,awr\_weights.py\}}, run against the
\texttt{mixed\_large} tier (50 real expert $+$ 100 \texttt{WeakPolicy}-rollout episodes, $78\%$
of frames from the wandering, task-\emph{never}-succeeding policy).
\begin{itemize}[leftmargin=1.4em,itemsep=2pt]
\item \emph{``The privileged-state oracle is bogus for a contact-manipulation task.''} Graded
against real expert step-gaps (a genuine dynamical ground truth on near-optimal data) rather
than assumed correct: Spearman $=+0.90$. It is fine.
\item \emph{``The latent graph itself is bad on Push-T.''} It is not --- Spearman
$=\mathbf{+0.9952}$ against expert step-gaps, better than Two-Room's own headline number
(\cref{sec:b0-results}). Median \texttt{phi\_dist} matches the true step-gap almost exactly
across scales: gaps $1/2/5/10/25/50/100 \to$ \texttt{phi\_dist} $1.0/2.0/5.0/10.0/25.0/50.0/99.0$.
\item \emph{``AWR's importance weights saturate into uniform behaviour cloning.''} Only $4.4\%$
of weights hit the clip (\texttt{awr\_weight\_clip=100}) --- not the cause either.
\end{itemize}

\textbf{The real cause.} \texttt{pseudo\_v} (the auxiliary regression target) is derived from
graph geodesic distance, which measures \emph{how many steps the data actually took} --- a
behavioural distance. IQL's expectile-$0.9$ value function is trained to learn the
\emph{optimal} cost-to-go. The two coincide on near-optimal data (exactly why the mechanism
works cleanly on Two-Room and on Push-T's own expert-only tiers) and diverge sharply once the
data contains wandering, never-succeeding trajectories. Measured directly on
\texttt{mixed\_large}'s own HER pairs: median \texttt{pseudo\_v} $=-5.85$ vs.\ the baseline's
TD-learned $V$ median $=-2.76$; auxphi's own $V$ gets dragged to $-4.98$. The downstream damage
is concrete: the AWR advantage picks up a systematic negative bias (mean $-0.24$ for auxphi vs.\
$-0.69$ for baseline), which \emph{inverts the actor's data preference} --- the ratio of mean
importance weight on expert vs.\ noisy transitions is $1.03$ for baseline (neutral) but only
$\mathbf{0.35}$ for auxphi, i.e.\ auxphi weights the never-succeeding \texttt{WeakPolicy} data
$\approx 3\times$ \emph{more} than expert data. This is also, run in reverse, exactly why the
mechanism works on Two-Room: its ``noisy'' tier is a noised \emph{expert} policy that still
solves the task, so behavioural and optimal cost-to-go never meaningfully diverge there.

\textbf{Data composition is the dominant lever --- more than any algorithmic knob tried.}

\begin{table}[h]
\centering
\small
\caption{Baseline live-rollout success rate by data composition (paper protocol,
\texttt{eval\_protocol=same\_episode}). Paper's own Push-T numbers, same protocol family:
GCBC~$75\%$, GCIVL~$33\%$, GCIQL~$20\%$, Random~$2\%$ (measured $4$--$9\%$ here on a comparable
pool).}
\begin{tabular}{llc}
\toprule
tier & composition & baseline success \\
\midrule
mixed\_large & 50 expert $+$ 100 WeakPolicy eps ($78\%$ wandering) & $26\%$ \\
expert\_1000 & 1000 pure expert eps & $\mathbf{54\%}$ \\
\bottomrule
\end{tabular}
\end{table}

A pure-expert-data baseline ($54\%$) already beats both of the paper's offline-RL baselines
(GCIQL $20\%$, GCIVL $33\%$) under their own protocol --- ``our method underperforms on
Push-T'' is not a supportable claim once tier composition is held fixed and stated.

\textbf{Two config-gated fixes, both defaulting to \texttt{false}} so no previously-published
Two-Room number moves (\texttt{config/graph/actor.yaml}): \texttt{awr\_normalize\_adv}
($z$-score the AWR advantage before $\exp(\cdot)$, so a biased/drifting $V$ cannot silently
rebias every actor weight) and \texttt{aux\_standardize} (regress $V$ onto the $z$-scored
\texttt{pseudo\_v} instead of its raw value, learning the graph's \emph{ordering} without
importing its behavioural absolute scale).

\begin{table}[h]
\centering
\small
\caption{\texttt{mixed\_large} auxphi live-rollout success under each fix in isolation
(\texttt{env=pusht tier=mixed\_large variant=auxphi}; baseline reference $26\%$, from above).}
\begin{tabular}{lc}
\toprule
condition & success rate \\
\midrule
auxphi (vanilla, both knobs \texttt{false}) & $14\%$ \\
auxphi $+$ \texttt{aux\_standardize=true} & $16\%$ \\
auxphi $+$ \texttt{awr\_normalize\_adv=true} & $\mathbf{21\%}$ \\
\bottomrule
\end{tabular}
\end{table}

\textbf{Takeaways.}
\begin{itemize}[leftmargin=1.4em,itemsep=2pt]
\item Never quote a Push-T policy number without stating tier composition --- expert-only vs.\
mixed changes baseline success by $2\times$, more than any method difference measured here.
\item Auxphi's benefit is conditional on the data being near-optimal; on a new environment it
should be tested on a pure-expert tier first, since a mixed/degraded tier can make it look like
an outright regression when the real cause is a distance-scale mismatch, not a broken mechanism.
\item The graph itself is the best-validated component in this whole pipeline on Push-T
($+0.9952$ vs.\ true step-gaps) --- it was never the suspect.
\item \texttt{awr\_normalize\_adv} recovers more of the gap than \texttt{aux\_standardize}
($21\%$ vs.\ $16\%$, vanilla $14\%$), suggesting the AWR-advantage bias is the larger of the two
mechanisms by which the behavioural/optimal mismatch does damage, though neither fix fully
closes it back to a clear win over the $26\%$ baseline.
\end{itemize}

\subsection{Within-episode identification edges}
\label{sec:b5-within-ep}

\textbf{Change.} \texttt{build\_id\_edges\_faiss\_capped} (the default, ``capped''
identification-edge construction used by Push-T) used to discard \emph{every} same-episode
$k$NN neighbour outright, regardless of how far apart in time the two frames were. Fixed to
accept a \texttt{min\_step\_gap} parameter (default $5$): a same-episode neighbour is now kept
as an identification edge as long as its step gap is $\ge$ \texttt{min\_step\_gap} --- a small
exclusion window avoids duplicating transition edges (gap $=1$) or near-duplicate consecutive
frames, while letting a wandering trajectory that loops back near an earlier state get the
shortcut its own revisit actually earns. Two-Room's \texttt{bruteforce} identification-edge path
is untouched, to protect its bit-exact historical numbers (\cref{sec:method}).

\textbf{Result} (gate check, Push-T \texttt{mixed} tier: 50 expert $+$ 50 \texttt{WeakPolicy}
episodes, $19{,}017$ landmarks, $300$k random eval pairs, oracle $=$ $z$-scored Euclidean over
projected state):

\begin{table}[h]
\centering
\small
\caption{Effect of allowing same-episode identification edges (\texttt{min\_step\_gap=5}).}
\begin{tabular}{lccc}
\toprule
construction & id.\ edges & same-episode Spearman ($n{=}3379$) & cross-episode Spearman ($n{=}296{,}595$) \\
\midrule
OLD (same-episode excluded) & $3{,}227$ (0 same-ep) & $0.560$ & $-0.178$ \\
NEW (\texttt{min\_step\_gap=5}) & $13{,}835$ (10,608 same-ep) & $\mathbf{0.678}$ & $-0.168$ \\
\bottomrule
\end{tabular}
\end{table}

Confirms the intended, targeted effect: $+10{,}608$ legitimate same-episode loop-back shortcuts
and $+0.12$ Spearman specifically on same-episode pairs, where the fix should matter.
Cross-episode/overall correlation is negative both before and after --- a separate,
already-diagnosed issue (\cref{sec:b5-calib-diagnosis,sec:b5-underperf-diagnosis}), not
something this fix targets.

\textbf{Live-rollout check.} Full actor training with the fix in place
(\texttt{scripts/actor\_train.py env=pusht tier=mixed variant=auxphi seed=0
eval\_protocol=cross\_episode}, $20$k steps, local GPU): peak held-out value-Spearman $0.375$ at
step $200$. \texttt{actor\_rollout\_eval.py ... select=rho n\_episodes=100} under the same
\texttt{cross\_episode} protocol (fixed held-out $(s,g)$ pairs, $250$-step budget):
$\mathbf{0/100}$ --- matching the pre-existing finding (\cref{sec:b5-b3-results}) that every
Push-T actor scores $0$ under this protocol.

\textbf{Takeaways.}
\begin{itemize}[leftmargin=1.4em,itemsep=2pt]
\item A real, targeted improvement to identification-edge quality on wandering data, verified at
the graph-construction level ($+0.12$ same-episode Spearman) --- but not a fix for the separate
\texttt{cross\_episode}-protocol $0\%$ finding: arbitrary cross-trajectory goals for a
contact-manipulation task, at a $250$-step budget, are just hard regardless of graph quality.
\item This fix is now baked into the default Push-T construction path (no config flag) --- every
Push-T result reported after this point, including \cref{sec:b5-stitching-eval}'s checkpoints,
uses it.
\end{itemize}

\subsection{Directed transition edges and a cross-episode auxiliary term}
\label{sec:b5-directed-cross-aux}

Two further changes to \texttt{scripts/common/graph\_lib.py} / \texttt{actor\_train.py}, tested
on Push-T's \texttt{mixed} tier.

\textbf{1. Directed transition edges.} \texttt{build\_weighted\_graph} used to add every
transition edge in both directions with equal weight. An env step is not generally reversible
(pushing a block cannot be undone by reversing the action), so an undirected geodesic overstated
closeness in the direction that matters. Transition edges are now one-directional (earlier-step
$\to$ later-step only); identification edges stay symmetric. Every \texttt{dijkstra(...)} call
across the active pipeline was switched to \texttt{directed=True} accordingly. \textbf{This also
changes Two-Room's historical bit-exact B0--B4 numbers} (\cref{sec:b0-results}--\cref{sec:b4ext-results}
no longer reproduce to full precision) --- accepted deliberately, since Two-Room's transitions
are only ``roughly'' reversible and preserving exact reproduction was never a design goal once a
genuinely more accurate graph is available. The six historical re-export shims and
\texttt{scripts/investigations/*} still call \texttt{dijkstra(..., directed=False)} and so
silently keep the old symmetric semantics --- correct, since those scripts are explicitly out of
scope for edits.

\textbf{2. A cross-episode auxiliary term.} The existing auxiliary-regression target
(\cref{alg:consumption}) only ever reads $\Phi$ at HER pairs, which are same-episode by
construction --- so $V$ never saw the graph's cross-episode stitching during training, even
though that is the graph's entire reason for existing over raw Euclidean distance. Fix: a
second, independently-sampled aux term at pairs drawn from the graph's own identification edges
that cross an episode boundary (real stitching links the graph itself found, not arbitrary
far-apart pairs, so $\Phi$ stays small by construction --- $3{,}860$ valid cross-episode training
pairs on the \texttt{mixed} tier, mean $\Phi=1.00$, i.e.\ mostly single direct edges). Weighted
by a new \texttt{aux\_cross\_lambda} (default $0.0$ $=$ off, exact historical behaviour).

\begin{table}[h]
\centering
\small
\caption{\texttt{mixed} tier, seed 0, \texttt{eval\_protocol=cross\_episode}, live-rollout
$n=100$ (Ada gnode079).}
\begin{tabular}{p{8.8cm}cc}
\toprule
condition & peak value-Spearman & live-rollout success \\
\midrule
historical (undirected graph, no within-ep id.\ edges, no cross aux) & $0.4219$ & $0/100$ \\
directed graph $+$ within-ep id.\ edges, \texttt{aux\_cross\_lambda=0} & $0.3764$@250 & $0/100$ \\
$+$ \texttt{aux\_cross\_lambda=0.3} & $0.3544$@200 & $0/100$ \\
\bottomrule
\end{tabular}
\end{table}

\textbf{Takeaways.}
\begin{itemize}[leftmargin=1.4em,itemsep=2pt]
\item Neither change flips the \texttt{cross\_episode}-protocol result in this single-seed test.
This is consistent with, not contrary to, \cref{sec:b5-within-ep}: every known Push-T condition
is floored at $0\%$ under this specific protocol, so it cannot show an improvement even if the
underlying mechanism works.
\item The peak-Spearman drop ($0.42\to0.38\to0.35$) is not apples-to-apples and should not be
read as these changes hurting --- directed transitions change what $\Phi$ measures (a harder,
more accurate directional distance), not a comparable-scale regression.
\item The more informative next test is \texttt{eval\_protocol=same\_episode} (where baseline
already scores above floor, \cref{sec:b5-eval-protocol,sec:b5-underperf-diagnosis}) with both
fixes in place --- the checkpoints used in \cref{sec:b5-stitching-eval} below were trained after
both landed (they are now unconditional defaults, not opt-in flags, except
\texttt{aux\_cross\_lambda}, which stays off there too).
\end{itemize}

\subsection{A stitching-specific showcase, with safeguards against p-hacking}
\label{sec:b5-stitching-eval}

\textbf{Motivation.} An arbitrary cross-episode $(s,g)$ pair is usually just \emph{far} in a
$250$-step-budget sense, regardless of whether the graph could find a shortcut --- so a plain
cross-episode eval is dominated by raw long-horizon difficulty, not by whether the graph can
stitch. A first naive attempt confirmed this: $100$ arbitrary cross-episode pairs, identical for
both actors, scored baseline $8\%$ / auxphi $9\%$ --- both near the random floor, telling us
nothing about stitching specifically. The actual claim under test --- that auxphi's edge over
baseline should be concentrated on cross-episode pairs the graph itself says are close via
identification edges (true stitching opportunities), and should vanish on pairs it says are far
(a negative control) --- needs pairs stratified by graph distance, not sampled arbitrarily.

\textbf{Methodology, with five safeguards fixed before looking at any outcome} (to guard against
post-hoc bucket-tuning), implemented in
\texttt{scripts/investigations/pusht\_lowscore/stitching\_eval.py}:
\begin{enumerate}[leftmargin=1.4em,itemsep=2pt]
\item Bucket boundaries fixed in advance and never tuned: near $<20$, mid $[20,50)$, far $\ge
50$, in true graph-distance (hop-count) units.
\item Every bucket reported, including far as the negative control.
\item Success is the real physical \texttt{eval\_state()} criterion, never $\Phi$ itself ---
$\Phi$ only decides which pairs to ask about.
\item The near bucket is independently cross-checked against the \emph{true} oracle distance
(\texttt{mech.build\_true\_distance\_oracle}, privileged $z$-scored state distance, computed
with no reference to the graph) --- rules out ``the graph invented a shortcut that isn't real.''
\item The near bucket is also reported by hop count (since every edge has weight $1$, $\Phi$
\emph{is} the shortest-path hop count) and raw latent Euclidean distance --- separates genuine
multi-hop stitching from a trivial near-duplicate pair findable by plain nearest-neighbour
lookup.
\end{enumerate}
Run as \texttt{python scripts/investigations/pusht\_lowscore/stitching\_eval.py env=pusht
tier=<expert\_1000|mixed\_large>}: $6{,}000$ cross-episode candidate pairs sampled uniformly at
random (seed $12345$) from the tier's own landmark set, bucketed by bounded-Dijkstra $\Phi$
(\texttt{limit=100}, \texttt{phi\_mode=sparse}); $40$ identical pairs per bucket run through
\emph{both} the baseline and auxphi checkpoints for that tier (same seed, same env, budget
$=250$ steps $=$ \texttt{mech.max\_episode\_steps}, the full budget rather than the paper's
$50$-step same-episode budget, since these are deliberately far cross-episode goals).
Checkpoints are the vanilla-auxphi ones trained after the fixes in
\cref{sec:b5-within-ep,sec:b5-directed-cross-aux} landed as defaults
(\texttt{aux\_cross\_lambda=0}, \texttt{aux\_standardize=false}, \texttt{awr\_normalize\_adv=false}),
selected by \texttt{select=rho}.

\textbf{Results.}

\begin{table}[h]
\centering
\small
\caption{Safeguards 4/5: near-bucket validation against the independent true oracle and
hop-count/Euclidean controls.}
\resizebox{\textwidth}{!}{%
\begin{tabular}{lcccccc}
\toprule
tier & near-bucket $n$ & oracle\_dist median & Spearman($\Phi$, oracle) & hop-count median & Euclidean median & \% trivial ($\Phi\le2$) \\
\midrule
expert\_1000 & 50/6000 & $0.823$ & $+0.4685$ & $14.0$ & $16.29$ & $0.0\%$ \\
mixed\_large & 499/6000 & $2.651$ & $+0.2833$ & $15.0$ & $12.34$ & $0.0\%$ \\
\bottomrule
\end{tabular}%
}
\end{table}

Both tiers' near buckets pass: genuinely multi-hop (median $14$--$15$ real graph hops, not a
single identification-edge shortcut), oracle-close (small median true state distance,
positively correlated with $\Phi$), and non-trivial ($0\%$ within $2$ hops of a plain
nearest-neighbour match).

\begin{table}[h]
\centering
\small
\caption{Live-rollout success by bucket, identical pairs for both variants, $40$/bucket,
$250$-step budget.}
\begin{tabular}{llcccc}
\toprule
tier & bucket & $\Phi$ range & baseline & \textbf{auxphi} & $\Delta$ \\
\midrule
expert\_1000 & near & $[0,20)$ & $\mathbf{45.0\%}$ & $37.5\%$ & $-7.5\%$ \\
expert\_1000 & mid  & $[20,50)$ & $7.5\%$ & $7.5\%$ & $+0.0\%$ \\
expert\_1000 & far  & $[50,\infty)$ & $2.5\%$ & $2.5\%$ & $+0.0\%$ \\
\midrule
mixed\_large & near & $[0,20)$ & $0.0\%$ & $0.0\%$ & $+0.0\%$ \\
mixed\_large & mid  & $[20,50)$ & $0.0\%$ & $0.0\%$ & $+0.0\%$ \\
mixed\_large & far  & $[50,\infty)$ & $0.0\%$ & $0.0\%$ & $+0.0\%$ \\
\bottomrule
\end{tabular}
\end{table}

This is a genuine negative result for the stitching-showcase hypothesis, reported as measured
rather than reframed: on \texttt{expert\_1000}, where both actors clear the floor at all
(near-bucket baseline $45\%$), \emph{baseline outperforms auxphi} by $7.5$ points on exactly the
pairs the graph says are the best stitching opportunities --- the opposite of the hypothesis.
Mid and far buckets are statistically tied and near the random floor for both variants, as
expected for a negative control. On \texttt{mixed\_large}, every bucket floors at $0\%$ for both
variants, so the eval has no discriminating power there at all (consistent with
\cref{sec:b5-underperf-diagnosis}'s finding that \texttt{mixed\_large}'s wandering-data
composition already suppresses auxphi's advantage broadly, not just on stitching pairs
specifically).

\textbf{Takeaways.}
\begin{itemize}[leftmargin=1.4em,itemsep=2pt]
\item The near buckets are methodologically sound (safeguards 4/5 pass on both tiers), so the
null/negative result is not an artifact of accidentally picking trivial or mislabelled pairs.
\item The result is consistent with \cref{sec:b5-underperf-diagnosis}'s mechanism: the auxiliary
regression target is a \emph{behavioural} distance (how many steps the data took along a path
the graph found), not the actor's own optimal cost-to-go. A genuine cross-episode stitching pair
is exactly where this mismatch should matter most, since the data itself never took that
shortcut --- so it is a plausible, not yet independently confirmed, explanation for why baseline
(which never queries the graph at cross-episode pairs at all) does no worse, and here slightly
better, than auxphi on the near bucket.
\item The as-yet-untried \texttt{aux\_cross\_lambda}$>0$ (\cref{sec:b5-directed-cross-aux}) is
the one mechanism in the current pipeline specifically designed to expose $V$ to the graph's
cross-episode opinion during training rather than only same-episode HER pairs --- a direct,
motivated next experiment for this eval, not yet run.
\item \texttt{expert\_1000}'s baseline being the strongest performer in the whole table ($45\%$
near-bucket) is itself informative: with enough pure-expert data, plain goal-conditioned BC/TD
needs no graph at all to find real cross-episode shortcuts through the data manifold --- the
mechanism, if it helps at all, would need to earn its keep against a substantially stronger, not
weaker, baseline than the small-data tiers used elsewhere in this document.
\end{itemize}

\subsection{A looser, empirically-set identification threshold, swept across data scale}
\label{sec:b5-eps99-sweep}

\textbf{Why a 99th-percentile threshold.} \Cref{sec:b5-calib-diagnosis} showed that
$\varepsilon^2$ (however it is set) is trying to separate two populations whose latent
distances barely overlap: real same-episode 1-step transitions cluster near $\varepsilon^2
\approx 1$--$2$, unrelated cross-episode pairs cluster near $\varepsilon^2 \approx 386$, four
orders of magnitude apart. Directly plotting both distributions confirms the calibrated
threshold ($\varepsilon^2{\approx}1.3$) sits at only the $50$th percentile of real transitions
--- it rejects half of all genuine 1-step moves as ``not similar enough,'' which is a far
stricter bar than ``is this a plausible identification.'' A threshold set at the
\emph{empirical} $99$th percentile of real transition distances is a direct, model-free way to
ask ``how loose can the threshold get while still keeping it inside the transition
population'': measured on the $300$-episode reference set, this is $\varepsilon^2 = 14.16$
--- notably not the same number \texttt{compute\_calibration}'s own chi2-formula gives at
$q=0.99$ ($3.12$; the real distribution has heavier tails than the chi2 approximation assumes,
a $4.5\times$ gap). At $\varepsilon^2=14.16$, only $0.016\%$ of random cross-episode pairs
would be accepted, so this is still a highly selective threshold, just $\sim\!10\times$ looser
than the calibrated default.

\begin{algorithm}[H]
\caption{Methodology: retraining setup and the $99$th-percentile threshold experiment}
\begin{algorithmic}[1]
\State \textbf{Retraining changes} (all tiers/variants below): $\gamma_{\text{HER}}: 0.9 \to
0.99$ (mean HER goal offset $10 \to 100$ steps, matching the paper's own GCIQL critic value);
$n_{\text{steps}}: 20000 \to 35000$; rollout-check episodes $12 \to 20$; success-checkpoint
tie-breaking changed from strict $>$ to $\ge$ (prefers the later, more-trained checkpoint on a
tie); only \texttt{select=success} evaluated going forward (\texttt{select=rho} tracks a
metric already shown, \cref{sec:b5-calib-diagnosis}, to be dominated by cross-episode pairs the
value function never trains on)
\State \textbf{Identification threshold:} \texttt{id\_eps2\_override=14.16} (new
\texttt{common.graph\_lib} parameter, bypasses \texttt{compute\_calibration} entirely for a
fixed, empirically-set $\varepsilon^2$ -- \cref{sec:b5-eps99-sweep} above)
\State \textbf{Tiers:} expert\_$\{100, 300, 1000\}$, both variants, 3 seeds
\State \textbf{Eval:} live-rollout success rate, 100 episodes, at three horizons --
goal$=\{25,50,100\}$ steps with a matching $2\times$ budget ($\{50,100,200\}$ steps)
\end{algorithmic}
\end{algorithm}

\textbf{Results.}

\begin{table}[h]
\centering
\small
\caption{Live-rollout success rate at $\varepsilon^2=14.16$ (empirical $99$th-percentile
threshold), mean over 3 seeds except where noted.}
\begin{tabular}{lccc}
\toprule
tier / goal & baseline & \textbf{auxphi} & winner \\
\midrule
expert\_100 / 25  & $21.7\% \pm 2.1$ & $18.3\% \pm 2.1$ & baseline \\
expert\_100 / 50  & $11.0\% \pm 1.4$ & $5.7\% \pm 2.1$  & baseline \\
expert\_100 / 100 & $2.0\% \pm 1.4$  & $0.7\% \pm 0.9$  & baseline \\
\addlinespace
expert\_300 / 25  & $31.7\% \pm 1.7$ & $33.3\% \pm 5.2$ & tie \\
expert\_300 / 50  & $13.3\% \pm 2.6$ & $12.7\% \pm 3.4$ & tie \\
expert\_300 / 100 & $8.3\% \pm 1.9$  & $5.7\% \pm 1.7$  & baseline \\
\addlinespace
expert\_1000 / 25  & $35.3\% \pm 4.6$ & $\mathbf{43.7\%} \pm 2.5$ & \textbf{auxphi} \\
expert\_1000 / 50  & $19.7\% \pm 4.2$ & $\mathbf{22.7\%} \pm 1.7$ & \textbf{auxphi} \\
expert\_1000 / 100 & $6.7\% \pm 2.6$  & $\mathbf{9.0\%} \pm 0.8$  & \textbf{auxphi} \\
\bottomrule
\end{tabular}
\end{table}

All rows above are the complete $n{=}3$ result. expert\_300's baseline row initially stalled at
seed 2 while building the dense $N{\times}N$ phi-distance matrix and had to be re-launched ---
root cause below.

\textbf{A clean scale trend.} Auxphi loses at expert\_100 (all 3 horizons), is statistically
tied at expert\_300, and wins outright at expert\_1000 (all 3 horizons, by $6$--$8$ points
absolute). This is the same explanation already proposed on other grounds
(\cref{sec:b5-calib-diagnosis}'s same-episode stitching-benefit comparison): more episodes
means more chances for a real cross-episode revisit to exist at all, even at a fixed threshold
-- so the auxiliary signal has more genuine substance to work with as data scales, independent
of anything about the threshold value itself. This sweep is the first place that trend shows up
directly on live-rollout success rather than a same-episode Spearman proxy.

\textbf{A real inefficiency, caught and fixed.} The stalled expert\_300/baseline/seed-2 run
traced to \texttt{actor\_train.py}'s tier setup (\texttt{get\_setup}/\texttt{build\_tier\_setups})
always building the dense phi-distance matrix by default, regardless of which variant will use
it -- baseline sets \texttt{aux\_lambda=0} and never reads \texttt{phi\_dist} at all, but
nothing skipped the (expensive, $O(n^2)$-memory) construction for it. With three combos per
tier building the same multi-GB matrix concurrently, the baseline runs were pure waste that
happened to tip a shared-memory budget over the edge. The existing \texttt{need\_graph=false}
config flag (already used for very large expert-only tiers, \cref{sec:method}) fixes this when
passed explicitly for baseline runs; it was not passed in this sweep's launch script and is now
corrected there for future runs.

\textbf{The same sweep at the default calibrated threshold.} To check whether the scale trend
above is specific to the looser $\varepsilon^2=14.16$ threshold, the identical sweep (same
tiers, retraining changes, and eval horizons) was re-run at \texttt{compute\_calibration}'s own
default threshold instead ($\varepsilon^2 \approx 1.3$, \texttt{run\_tag=\_defaulteps}, no
override).

\begin{table}[h]
\centering
\small
\caption{Live-rollout success rate at the default calibrated threshold ($\varepsilon^2 \approx
1.3$), mean over 3 seeds.}
\begin{tabular}{lccc}
\toprule
tier / goal & baseline & \textbf{auxphi} & winner \\
\midrule
expert\_100 / 25  & $23.7\% \pm 3.9$ & $\mathbf{27.0\%} \pm 1.6$ & \textbf{auxphi} \\
expert\_100 / 50  & $\mathbf{9.7\%} \pm 2.9$  & $7.3\% \pm 2.5$  & baseline \\
expert\_100 / 100 & $1.7\% \pm 0.9$  & $1.3\% \pm 0.5$  & tie \\
\addlinespace
expert\_300 / 25  & $30.7\% \pm 1.2$ & $30.3\% \pm 2.6$ & tie \\
expert\_300 / 50  & $12.3\% \pm 1.7$ & $14.0\% \pm 2.2$ & tie \\
expert\_300 / 100 & $6.3\% \pm 0.9$  & $5.7\% \pm 1.7$  & tie \\
\addlinespace
expert\_1000 / 25  & $33.7\% \pm 4.2$ & $\mathbf{38.7\%} \pm 1.7$ & \textbf{auxphi} \\
expert\_1000 / 50  & $17.7\% \pm 1.9$ & $\mathbf{26.7\%} \pm 5.4$ & \textbf{auxphi} \\
expert\_1000 / 100 & $7.0\% \pm 0.8$  & $\mathbf{8.7\%} \pm 0.9$  & \textbf{auxphi} \\
\bottomrule
\end{tabular}
\end{table}

All rows above are the complete $n{=}3$ result. expert\_300/auxphi/seed\,1 initially crashed
reading a half-written phi-distance cache file (three auxphi seeds for the same tier tried to
build and cache that matrix concurrently) and was retrained separately once the cache was
cleared.

\textbf{The scale trend replicates.} The same pattern holds at the default threshold: auxphi
is mixed-to-behind at expert\_100 (wins the shortest horizon, loses the other two), a
statistical tie at expert\_300, and a clean win at expert\_1000 across all three horizons. So
the earlier trend was not an artifact of the looser $\varepsilon^2=14.16$ threshold --
identification-edge density growing with data scale is what drives it, independent of exactly
where the threshold is set, consistent with \cref{sec:b5-calib-diagnosis}'s explanation that
more episodes means more genuine cross-episode revisits for the auxiliary signal to learn from.

\subsection{Testing the advantage-drift hypothesis: \texttt{aux\_standardize} and \texttt{awr\_normalize\_adv}}
\label{sec:b5-stdadv-sweep}

\textbf{Why this test.} A comment already in \texttt{actor.yaml}, from an earlier
investigation, records a measured advantage drift on Push-T: mean AWR advantage
$-0.689$ (baseline) vs.\ $-0.236$ (auxphi) on \texttt{mixed\_large}. Since AWR's policy
weight is $\exp(\beta \cdot \text{advantage})$ clamped at \texttt{awr\_weight\_clip}, a
uniform shift in advantage is not cosmetic: it changes how many actions saturate at the
clamp, which flattens exactly the fine-grained preference ranking AWR depends on. Two
existing (never-enabled) config flags target this directly:
\texttt{aux\_standardize} regresses $V$ onto the graph pseudo-value's z-scored
\emph{ordering} rather than its raw behavioral-distance scale; \texttt{awr\_normalize\_adv}
z-scores the advantage itself immediately before the $\exp(\cdot)$ weighting, so a
constant drift from the aux term cannot silently change the clamp population. Neither
was enabled in any Push-T run before this one, including \cref{sec:b5-eps99-sweep}'s
two sweeps.

\begin{algorithm}[H]
\caption{Methodology: identical to \cref{sec:b5-eps99-sweep}'s default-threshold sweep, plus the two flags}
\begin{algorithmic}[1]
\State \textbf{Unchanged from \texttt{run\_defaulteps\_sweep.sh}:} $\varepsilon^2 \approx 1.3$ (\texttt{compute\_calibration} default, no override); $\gamma_{\text{HER}}=0.99$; $n_{\text{steps}}=35000$; rollout-check episodes $=20$; tiers expert\_$\{100,300,1000\}$, both variants, 3 seeds; eval at goal$=\{25,50,100\}$ with matching $2\times$ budget; \texttt{select=success} only
\State \textbf{Added:} \texttt{aux\_standardize=true awr\_normalize\_adv=true}
\State \textbf{Distinct \texttt{run\_tag=\_defaulteps\_stdadv}} so this cannot collide with or resume from the plain \texttt{\_defaulteps} checkpoints; same tier-cache directory reused (the graph/$\Phi$ matrix is unchanged, only the training loss is)
\end{algorithmic}
\end{algorithm}

\textbf{Results.}

\begin{table}[h]
\centering
\small
\caption{Live-rollout success rate at the default calibrated threshold, with \texttt{aux\_standardize}/\texttt{awr\_normalize\_adv} enabled, mean over 3 seeds.}
\begin{tabular}{lccc}
\toprule
tier / goal & baseline & auxphi & winner \\
\midrule
expert\_100 / 25  & $\mathbf{26.7\%} \pm 0.5$ & $22.7\% \pm 1.9$ & \textbf{baseline} \\
expert\_100 / 50  & $10.7\% \pm 2.1$ & $10.0\% \pm 0.8$ & tie \\
expert\_100 / 100 & $2.0\% \pm 0.8$  & $1.3\% \pm 1.2$  & tie \\
\addlinespace
expert\_300 / 25  & $38.3\% \pm 3.3$ & $34.0\% \pm 6.5$ & tie \\
expert\_300 / 50  & $\mathbf{18.3\%} \pm 2.4$ & $13.3\% \pm 3.3$ & \textbf{baseline} \\
expert\_300 / 100 & $4.3\% \pm 0.9$  & $6.3\% \pm 1.2$  & tie \\
\addlinespace
expert\_1000 / 25  & $\mathbf{48.0\%} \pm 3.6$ & $41.0\% \pm 5.9$ & \textbf{baseline} \\
expert\_1000 / 50  & $26.7\% \pm 3.3$ & $25.3\% \pm 3.9$ & tie \\
expert\_1000 / 100 & $7.7\% \pm 1.9$  & $10.7\% \pm 3.7$ & tie \\
\bottomrule
\end{tabular}
\end{table}

\textbf{The hypothesis does not hold up --- if anything, it runs backward.} Under the
plain \texttt{\_defaulteps} sweep (\cref{sec:b5-eps99-sweep}'s second table), auxphi had
$4$ clear wins (expert\_100/25, and all three expert\_1000 horizons by $5$--$9$ points) against
$1$ clear baseline win and $4$ ties. With the drift fix applied, auxphi has \emph{zero} clear
wins anywhere: baseline picks up $3$ clear wins (expert\_100/25, expert\_300/50,
expert\_1000/25) and everything else is now a tie. Most notably, expert\_1000's clean
$3$-for-$3$ auxphi sweep --- the strongest evidence for the graph signal anywhere in the
Push-T results, and the basis for \cref{sec:b5-eps99-sweep}'s scale-trend claim --- is gone:
baseline now leads or ties on all three of its horizons. Both baseline's and auxphi's
absolute numbers move under the fix (advantage normalization is not auxphi-specific, since
\texttt{awr\_normalize\_adv} applies to both variants), but baseline moves up more often and
by more, particularly at the larger tiers and shortest horizon (expert\_1000/25: $+14$
points for baseline vs.\ $+2$ for auxphi). The advantage-drift measured on
\texttt{mixed\_large} was real, but correcting it did not unlock a hidden auxphi advantage
--- it mostly stabilized baseline's own AWR extraction, which had nothing to do with the
aux term in the first place. This narrows, rather than resolves, the open question from
\cref{sec:b5-calib-diagnosis}: the scale-trend explanation (more episodes $\to$ more genuine
stitching substance) is not wrong on its own terms, but it was measured under a training
setup with an uncontrolled confound, and the one sweep that isolates that confound erases
the effect it was meant to explain.

\section{Risks}

\begin{enumerate}[leftmargin=1.4em,itemsep=3pt]
\item \textbf{Potential-based shaping (mode 3) fired, and not the way this paragraph predicted ---
but the underlying idea recovers under a different consumption mode.} The hypothesis here was
that shaping would help \emph{most} in degraded (scarce/mixed) data tiers, where stitching matters
most and TD has least to work with, with clean-data results serving as a control. The measured
result was the reverse: mode 3 won only with 100 clean episodes, lost at every smaller size, and
lost \emph{worse} as more heterogeneous data was added (\cref{sec:b3-results}). That was a
confirmed risk in the direction least favorable for deployment. What resolved it was not tuning
mode 3 (scaling its strength down closes most, not all, of the gap) but switching to mode 2 ---
using $\Phi$ as a direct auxiliary regression target instead of folding it into the TD-bootstrapped
reward (\cref{sec:method}). Mode 2 wins at 5 of 6 tiers, by a margin that grows with data. The
remaining risk is narrower: it is now about \emph{which consumption mode} to recommend by
default, not about whether the graph-distance signal is worth anything --- mode 3's
policy-invariance guarantee is real but does not, by itself, protect against noise compounding
through TD.
\item \textbf{The graph is dense, and it gets denser faster than the data does.} Measured, not
estimated: identification edges scale as landmark-count$^{1.96}$ on Two-Room
(\cref{sec:b1-results}), extrapolating to $\sim 1.5\times 10^{10}$ edges at Push-T's scale.
Plain landmark subsampling delays this but does not fix the scaling law. The construction needs
a top-$k$-per-node cap (bounding degree directly, turning quadratic growth into $O(nk)$) before
any Push-T attempt; this is now a blocking prerequisite for the Push-T work (\cref{sec:pusht}),
confirmed in B1 rather than assumed.
\item \textbf{Physically close, dynamically separated.} States on opposite sides of a thin wall
can be falsely identified --- the measured false-edge rate is 2--7\% depending on $q$.
Successor-consistency filtering (if $i,j$ are the same state under similar actions, their
recorded successors must also pass the test) was implemented and tested in B1: it removed 8.4\%
of identification edges but left the false-edge rate and downstream Spearman unchanged, because
only 25.6\% of edges had similar-enough actions to even test --- Two-Room's target is invisible
in the pixel observation, so same-position pairs from different episodes are typically executing
unrelated goal-directed policies. The false-edge problem itself is real and unresolved; this
particular fix does not solve it here, and should not be assumed to transfer to an environment
where the goal is visible.
\item \textbf{Novelty is narrower than it looks.} SoRB (Eysenbach et al.\ 2019) and semi-parametric
topological memory (Savinov et al.\ 2018) both build graphs over experience; both learn their edge
predicate from online or reward-labelled data. The specific claim here is narrower and should be
stated narrowly: \emph{a reward-free, frozen, pretrained latent metric with a
distributionally-calibrated identification threshold, deployed as a policy-invariant potential.}
B2 is what decides whether even that survives.
\end{enumerate}

\section{Resources}

Two-Room data and the pretrained checkpoint are downloaded and working; B0's encoded landmarks
are cached (\texttt{outputs/b0\_tworoom/landmarks\_ep300.npz}) and reusable without re-encoding.
Push-T data is already local. Baseline checkpoints (PLDM, LeJEPA, DINO-WM, GCIQL, GCBC) are
published for Two-Room, so B2 and B3 need no baseline training from scratch. Encoding runs at
$\sim$600 frames/s on the local GPU; B0 end-to-end took under two minutes with landmarks cached.

The loader (\texttt{scripts/tworoom\_lewm\_loader.py}) contains a state-dict key remap needed
because the installed \texttt{transformers} renamed ViT block internals relative to the
published checkpoint; verified as a pure rename before use. Any other pretrained LeWM checkpoint
will need the same treatment.

\end{document}
